%% 物理第2巻『熱学』：組 熱W1（第22回編集委員会の手引き tools/BRIEF-phys2.md）
%% p2-02.tex  熱-2　仕事で温まる：第一法則と内部エネルギー（★2〜4。3段にしない）
%% 第一法則の符号（手引き2章の3・計画書 A-7、この巻の決まりをこの章で決める）：
%%   本文は ΔU＝Q＋W。Q は系が「受けとる」熱、W は系が「される」仕事（高校の教科書・IUPAC と同じ）。
%%   系が「する」仕事は W′＝−W と書く（数学 8-7⑩ の「気体が外にする仕事 ∫p dV」は W′）。
%%   紙面では矢印の図で、入る熱 Q_入・出る熱 Q_出・される仕事 W_入・する仕事 W_出 をどれも正の数で書き、ΔU＝Q_入−Q_出＋W_入−W_出 と立てる。
%%   機械の本の形 Q＝ΔU＋W（W は系がした仕事）は「せいかくに言うと」と「専門書をのぞいてみよう」で見せる。
%%   サイクルの章（熱-5・7・8）は W_出・W_入 の添字で分けてよい（記号カードに書いた）。
%% 記号（この章）：内部エネルギー U〔J〕、熱量 Q、される仕事 W、する仕事 W′、小さな熱・仕事 δQ・δW（数学 8-3 の専門書と同じ）、
%%   圧力 p〔Pa〕、体積 V、絶対温度 T〔K〕・セルシウス温度 θ〔℃〕（熱-1 ②）、仕事率・電力 P〔W〕（圧力 p とは大文字・小文字で分ける）、
%%   運動エネルギー K（力-4 と同じ）、重さは mg（W と書かない）、高さ h、比熱 c、熱容量 C_th。
%% 数値はすべて Python で計算した：checks/物理第2巻/熱-2/check_p2_02.py（出力 check_p2_02_out.txt）
%% 第1巻・数学の本との約束：
%%   力-4 ⑦「熱まで入れたエネルギーの数え方（熱力学の第一法則）は、熱-2 で学びます」→ ④の「力学的エネルギーと内部エネルギー」・例題6（ブレーキ）・例題7（自転車で下る）。
%%   数学 8-7⑩「物理の熱力学（熱-2）では…仕事や熱は道によって変わるが、内部エネルギーは状態だけで決まる」→ ③④・例題5（同じ3つの道 400・200・300 J）。
%%   数学 8-3 の専門書「dU は完全微分であるが、熱 δQ と仕事 δW は完全微分ではない（物理 熱-2）」→ ⑦・例題11。
%%   数学 8-3 の物理での使い方「定積比熱と定圧比熱（物理 熱-2・熱-10）」→ ⑦の終わり（くわしくは熱-3 ④）。
%%   空気入れの数（20 ℃・0.20 L・体積1/3・仕事 28 J・約182 ℃）は、熱-3 ⑥（例題6）と同じ題材・同じ数にした（計画書 C-4 の2）。
%% 出典（本文で使った事実・値）：
%%   ・ランフォード（ベンジャミン・トンプソン、1753〜1814年）：1798年、大砲の砲身に穴をあける作業で、摩擦の熱が尽きないことを報告（Phil. Trans. R. Soc. 88, 80 (1798)）。
%%   ・マイヤー 1842年（熱と仕事の換算を気体の比熱から見積もった）。
%%   ・ジュール（1818〜1889年）：羽根車の実験は1845年から。1850年の論文 "On the Mechanical Equivalent of Heat"（Phil. Trans. R. Soc. 140, 61–82）で
%%     「水1ポンドを華氏1度上げる熱＝772ポンドの物を1フィート持ち上げる仕事」（772.692 ft·lbf）。J/cal に直すと約 4.16。
%%   ・ジュールの自由膨張の実験：1845年（Phil. Mag. 26, 369）。約22気圧の空気と真空の2つの銅の容器を水の中でつなぎ、水の温度は変わらなかった。
%%     ゲイ＝リュサックも1807年ごろに似た実験をした。精密な実験では、ふつうの気体はわずかに温度が下がる（理想気体からのずれ。熱-10）。
%%   ・第一法則の符号：高校の教科書は ΔU＝Q＋W（された仕事）の形が多い〔要確認：各社の教科書の本文〕。IUPAC も ΔU＝q＋w（系にされた仕事）。
%%     機械工学の教科書は Q＝ΔU＋W（系がした仕事）の形（日本機械学会『JSME テキストシリーズ 熱力学』など）〔要確認：本文の文言〕。
%%   ・1 cal＝4.184 J（熱化学カロリー。熱-1 と同じ）。1日2000 kcal＝8.37 MJ。
%%   ・ブレーキの円板（鋳鉄）の比熱 約0.46 kJ/(kg·K)（ねずみ鋳鉄。〔要確認：JIS G 5501 の附属書・メーカーの資料〕）。
%%   ・空気の定積比熱 c_v＝718 J/(kg·K)、密度 1.20 kg/m^3（20 ℃・1気圧）：CoolProp 8.0 で計算（717.7、1.2046）。
%%   ・確率熱力学（ゆらぎの大きい小さな系の熱と仕事）：関本謙 K. Sekimoto, J. Phys. Soc. Jpn. 66, 1234 (1997)。
%%     実験の例：S. Toyabe et al., Nature Physics 6, 988 (2010)（情報を使って、熱ゆらぎから仕事を取り出す）。
\makeatletter
% ピストンとシリンダー（横向き）  \qxcpiston{(左下)}{長さ}{高さ}{ピストンの位置}
\newcommand\qxcpiston[4]{\begin{scope}[shift={#1}]
  \fill[bkBlue!15] (0,0) rectangle (#4,#3);
  \draw[bkInk, line width=0.9pt] (#2,0) -- (0,0) -- (0,#3) -- (#2,#3);
  \filldraw[fill=bkGray!45, draw=bkInk, line width=0.7pt] (#4,0.4) rectangle ({#4+3},{#3-0.4});
  \filldraw[fill=bkGray!45, draw=bkInk, line width=0.7pt] ({#4+3},{#3/2-0.8}) rectangle ({#2+6},{#3/2+0.8});
\end{scope}}
\makeatother
\tikzset{qxc lab/.style={font=\sffamily\footnotesize, inner sep=1pt},
  qxc arr/.style={-{Stealth[length=4.4pt, width=4pt]}, line width=1.6pt}}

\phchapter{2}{仕事で温まる：\\第一法則と内部エネルギー}{自転車の空気入れの筒の先は、なぜ熱くなる？}{\bklevelrange{2}{4}}{仕事で温まる：第一法則と内部エネルギー}
\bkmichishirube{3-7,5-12,8-2,8-3,8-7,力-4,熱-1}{熱-3,熱-4,熱-5,熱-6}

\begin{goal}
\begin{itemize}[label={\color{bkTheme}●}, leftmargin=1.6\zw, itemsep=0pt, topsep=1pt]
\item \bklevel{2}\enspace 熱と仕事が同じエネルギーであることを、実験から説明できる。
\item \bklevel{2}\enspace 第一法則 $\Delta U=Q+W$ で、エネルギーの出入りを数えられる。
\item \bklevel{3}\enspace $p$-$V$ 図の面積から、気体がする仕事を求められる。
\item \bklevel{4}\enspace 状態量と、道すじで変わる量を、全微分で見分けられる。
\end{itemize}
\end{goal}

\phkoushiki{%
\phk{熱と仕事は同じエネルギー（ジュールの実験、熱の仕事当量）}{高校物理基礎}{ph:熱-2:joule}
\phk{内部エネルギーと熱力学の第一法則 $\Delta U=Q+W$}{高校物理基礎}{ph:熱-2:dai1}
\phk{気体がする仕事 $W'=p\,\Delta V$（$p$-$V$ 図の面積）}{高校物理}{ph:熱-2:pdV}
\phk{断熱自由膨張では $\Delta U=0$（理想気体の温度は変わらない）}{高校物理}{ph:熱-2:jiyuu}
}

\begin{hatena}
\noindent\begin{minipage}[c]{0.58\linewidth}\raggedright
自転車のタイヤに、空気入れで空気を入れます。何十回もいきおいよく押したあとで筒をさわると、下のほう（ホースの付け根に近い、筒の先）が、はっきり温かくなっています。

火も電気も使っていないのに、なぜ温まるのでしょう。しかも、ピストンが擦れる筒の上のほうより、筒の先のほうが熱いのです。
\end{minipage}\hfill
\begin{minipage}[c]{0.38\linewidth}\centering
\begin{tikzpicture}[x=1mm, y=1mm]
  \filldraw[fill=bkGray!15, draw=bkInk, line width=0.8pt] (10,4) rectangle (20,44);
  \filldraw[fill=bkRed!35, draw=none] (10.4,4.4) rectangle (19.6,14);
  \draw[bkInk, line width=0.8pt] (10,4) rectangle (20,44);
  \filldraw[fill=bkGray!45, draw=bkInk, line width=0.7pt] (10.4,26) rectangle (19.6,29);
  \draw[bkInk, line width=1.2pt] (15,29) -- (15,50);
  \draw[bkInk, line width=2.2pt, rounded corners=1mm] (7,50) -- (23,50);
  \filldraw[fill=bkGray!40, draw=bkInk, line width=0.7pt] (2,0) rectangle (28,4);
  \draw[bkInk, line width=1.4pt] (20,7) .. controls (30,7) and (34,12) .. (36,22);
  \draw[qxc arr, bkBlue!80!black] (26,52) -- (26,40);
  \node[qxc lab, anchor=west, text=bkRed!75!black] at (21,11) {熱い};
  \node[qxc lab, anchor=south, text=bkBlue!80!black] at (26,52.5) {押す};
\end{tikzpicture}
\end{minipage}
\end{hatena}

\noindent 熱-1 では、熱がエネルギーで、J で数えることを使いました。この章では、そのことを確かめた実験から始め、\textbf{熱と仕事を合わせて数える}きまり、熱力学の第一法則を作ります。第1巻の力-4 で、摩擦がした負の仕事の分の力学的エネルギーは「熱や音になる」と書きました。そのゆくえを、ここで数で追いかけます。

\begin{yattemiyou}[擦ったり、押し縮めたりすると、温まる？]
\textbf{用意するもの}\par
自分の両手、自転車の空気入れとタイヤ（大人といっしょに）、時計\par\vspace{2pt}
\noindent\begin{minipage}[c]{0.62\linewidth}
\begin{enumerate}[label={\bkcircled[bkGreen]{\arabic*}}, leftmargin=2.4\zw, itemsep=1pt, topsep=0pt]
\item 両手のひらを、10秒間、強く擦り合わせる。すぐにほおに当てて、擦る前とくらべる。
\item 空気入れで、空気の少ないタイヤに、30回ほどいきおいよく空気を入れる。
\item すぐに、筒の上のほう（ピストンが擦れる所）と、下のほう（筒の先）を手でさわって、温かさをくらべる。
\end{enumerate}
\end{minipage}\hfill
\begin{minipage}[c]{0.34\linewidth}\centering
\begin{tikzpicture}[x=1mm, y=1mm]
  \filldraw[fill=bkOrange!25, draw=bkInk, line width=0.7pt, rounded corners=3mm] (0,0) rectangle (14,26);
  \filldraw[fill=bkOrange!25, draw=bkInk, line width=0.7pt, rounded corners=3mm] (16,0) rectangle (30,26);
  \draw[{Stealth[length=3pt]}-{Stealth[length=3pt]}, bkRed, line width=0.9pt] (8,30) -- (22,30);
  \node[qxc lab, anchor=south] at (15,31) {擦る};
\end{tikzpicture}
\end{minipage}\par\vspace{6pt}
{\small 空気入れの筒の先は、温かくなっても、やけどするほどにはなりません。タイヤには決められた圧力以上に入れないようにします。}
\end{yattemiyou}

\noindent やってみようの答え：擦った手は、温かくなります。手が手を押しながら動かす力（摩擦力）がした仕事が、皮ふの温度を上げたのです。空気入れでは、筒の先のほうが、上のほうより温かく感じます。擦れる所よりも、空気が強く押し縮められる所のほうが温まっています。なぜそうなるのかは、⑤で考えます。

\begin{omoidasou}
\bkrecall{力-4}{仕事 $W=\int_C\vec F\cdot\dd\vec r$。運動エネルギーの変化＝された仕事。摩擦の負の仕事の分、力学的エネルギーが減る。}
\bkrecall{5-12}{力が変わるときの仕事は $\int F\,\dd x$。グラフの下の面積。}
\bkrecall{3-7}{温度を変えずに押し縮めた注射器の空気は、$pV$ が一定（約 2000\,kPa·mL）。}
\bkrecall{8-7}{$p$-$V$ 図の道。1\,L・200\,kPa から 3\,L・100\,kPa への3つの道で、気体が外にする仕事 $\int p\,\dd V$ は 400\,J・200\,J・300\,J。}
\bkrecall{8-3}{全微分 $\dd U=\bigl(\partial U/\partial T\bigr)_V\dd T+\bigl(\partial U/\partial V\bigr)_T\dd V$。熱と仕事は全微分ではない印に $\delta$。}
\bkrecall{8-2}{2回の偏微分は、順番を変えても同じ（$\partial^2f/\partial x\,\partial y=\partial^2f/\partial y\,\partial x$）。}
\bkrecall{熱-1}{熱量 $Q$〔J〕。$Q=mc\,\Delta T$。熱のバランス式「たまる熱＝入る熱－出る熱」。}
\end{omoidasou}

%% ======================== 物理の道具 ========================
\bksection{doguu}{物理の道具}

\bksubsection{1}{熱の正体：熱素からエネルギーへ}
\label{ph:熱-2:sec1}
\phboxhead{熱は「物質」だと考えられていた}
18世紀の科学者の多くは、熱を、目に見えず重さもない「\textbf{熱素}」という物質だと考えていました。熱い物には熱素が多く、熱い物から冷たい物へ熱素が流れこむので温まる、という考えです。熱-1 で学んだ熱のバランス式（熱量の保存）も、熱素がなくならない物質だと考えると、うまく説明できました。

\phboxhead{擦ると、熱はいくらでも出てくる}
ところが、1798年、ランフォード（ベンジャミン・トンプソン）は、大砲の砲身に穴をあける工場で、ふしぎなことに気づきました。馬に機械を回させて、鈍い刃で金属をけずり続けると、熱がいつまでも出てきて、まわりの水がわくほどになったのです。熱素がもとから金属の中にあるだけなら、いつかは出つくすはずです。ランフォードは、「熱は物質ではなく、擦る動き（仕事）から生まれる」と考えました。

\phboxhead{ジュールの羽根車の実験}
\begin{bkfigure}
\begin{minipage}[c]{0.46\linewidth}\centering
\begin{tikzpicture}[x=0.72mm, y=0.8mm]
  \filldraw[fill=bkBlue!18, draw=bkInk, line width=0.8pt] (10,0) rectangle (38,24);
  \draw[bkInk, line width=0.9pt] (24,4) -- (24,36);
  \foreach \yy in {6,12,18} {\filldraw[fill=bkGray!45, draw=bkInk, line width=0.5pt] (15,\yy) rectangle (33,{\yy+2});}
  \filldraw[fill=bkGray!30, draw=bkInk, line width=0.6pt] (20,34) rectangle (28,38);
  \draw[bkInk, line width=0.6pt] (28,36) -- (52,36) (20,36) -- (-4,36);
  \filldraw[fill=bkGray!20, draw=bkInk, line width=0.6pt] (52,36) circle (2.4) (-4,36) circle (2.4);
  \draw[bkInk, line width=0.6pt] (54.4,36) -- (54.4,14) (-6.4,36) -- (-6.4,14);
  \filldraw[fill=bkOrange!45, draw=bkInk, line width=0.6pt] (51,6) rectangle (58,14) (-10,6) rectangle (-3,14);
  \draw[qxc arr, bkRed!80!black] (61,13) -- (61,3);
  \node[qxc lab, anchor=west, text=bkRed!75!black] at (62,8) {落ちる};
  \node[qxc lab, anchor=north] at (24,-1.5) {水の中の羽根車};
  \node[qxc lab, anchor=north] at (-6.5,5) {重り};
\end{tikzpicture}
\end{minipage}\hfill
\begin{minipage}[c]{0.5\linewidth}\raggedright
ジュールは1845年ごろから、重りが落ちる力で水の中の羽根車を回す実験を続けました。水は少しだけ温まります。落ちた高さから重力の仕事 $mgh$ が、水の温度の上がりから水が受けとった熱 $mc\,\Delta T$ がわかります。
\end{minipage}
\bkfigcaption{ジュールの実験（しくみを簡単にかいた）。重りがした仕事が、水をかき回して温める。}\label{ph:熱-2:fig-joule}
\end{bkfigure}

ジュールは、1850年の論文で、「水1ポンド（約\qty{454}{g}）を華氏1度上げる熱は、772ポンドの物を1フィート持ち上げる仕事と同じ」とまとめました。今の単位に直すと、1\,cal の熱が約\qty{4.16}{J}の仕事にあたります。この比を\textbf{熱の仕事当量}といいます。\bkidx{ねつのしごとうりょう}{熱の仕事当量}いまは cal を $\qty{1}{cal}=\qty{4.184}{J}$ と決めているので（熱-1 ③）、熱の仕事当量は単位の換算の数になりました。

\label{ph:熱-2:joule}こうして、\textbf{熱と仕事は、同じエネルギーの2つの移り方}だとわかりました。仕事は「力で押して動かす」移り方、熱は「温度の差で、分子どうしがぶつかって伝わる」移り方です。だから、どちらも J で数えて、たしたりひいたりできます。

\begin{phyomi}[熱の仕事当量 $\qty{1}{cal}=\qty{4.184}{J}$]
「水\qty{1}{g}を\qty{1}{K}温める熱は、約\qty{0.43}{kg}の物を\qty{1}{m}持ち上げる仕事と同じ」（$4.184\div9.8\approx0.43$）。熱にすると小さな量でも、仕事にすると大きい。
\end{phyomi}

\begin{phbaisuru}{重りの落ちる回数}
重り\qty{10}{kg}が\qty{2}{m}落ちると、重力の仕事は $10\times9.8\times2=\qty{196}{J}$ です。水\qty{1}{kg}の温度は、1回で約 $196\div4200\approx\qty{0.047}{K}$ しか上がりません。20回繰り返して約\qty{0.93}{K}、40回で約\qty{1.9}{K}です（例題1）。ジュールは、温度計を\qty{0.01}{K}より細かく読む工夫をして、この小さな変化を測りました。
\end{phbaisuru}

%% ---------------------------------------------------------------
\bksubsection{2}{系とまわり、状態量と内部エネルギー}
\label{ph:熱-2:sec2}
\phboxhead{系・まわり・境界}
熱や仕事を数えるときは、まず、どこからどこまでを考えるかを決めます。考える物のまとまりを\textbf{系}（けい）、そのほかを\textbf{まわり}、2つを分ける面を\textbf{境界}といいます。\bkidx{けい}{系}シリンダーの中の空気を系にするなら、シリンダーの壁とピストンの内側の面が境界です。この章では、境界を通って物（質量）が出入りしない系（\textbf{閉じた系}）を考えます。物が流れこんだり出ていったりする系は、熱-4 で考えます。

\begin{bkfigure}
\begin{tikzpicture}[x=1mm, y=1mm]
  \qxcpiston{(30,0)}{50}{24}{34}
  \draw[bkRed!80!black, line width=1pt, dash pattern=on 4pt off 2pt, rounded corners=1mm] (29,-1) rectangle (64.6,25);
  \node[qxc lab, text=bkBlue!80!black] at (46,12) {系（中の気体）};
  \node[qxc lab, anchor=south, text=bkRed!75!black] at (46,26) {境界};
  \draw[qxc arr, bkOrange!85!black] (10,12) -- (28,12);
  \node[qxc lab, anchor=south, text=bkOrange!80!black] at (16,13) {熱 $Q$};
  \draw[qxc arr, bkPurple!80!black] (100,12) -- (88,12);
  \node[qxc lab, anchor=south, text=bkPurple!80!black] at (98,13.5) {仕事 $W$};
  \node[qxc lab, anchor=north, text=bkGray] at (60,-3) {まわり（外の空気・ガスの火・押す手など）};
\end{tikzpicture}
\bkfigcaption{系・まわり・境界。熱と仕事は、境界を通って移るエネルギー。}\label{ph:熱-2:fig-kei}
\end{bkfigure}

\phboxhead{状態量}
系の、いまの状態だけで決まる量を\textbf{状態量}といいます。\bkidx{じょうたいりょう}{状態量}気体なら、圧力 $p$、体積 $V$、温度 $T$ が状態量です。どんなふうにして今の状態になったか（温めたのか、押し縮めたのか）には、よりません。数学 8-3①の、山の高さが「いる場所だけで決まる」のと同じです。

\phboxhead{内部エネルギー $U$}
系の中の分子は、いつも熱運動をしています（熱-1 ②）。分子1個1個の運動エネルギーと、分子どうしが引き合う力による位置エネルギーを、系全体でたしたものを\textbf{内部エネルギー}といい、$U$ で書きます。\bkidx{ないぶえねるぎー}{内部エネルギー}温度が上がると、分子の動きが速くなるので、$U$ は増えます。氷がとけるときは、温度は変わりませんが、分子のつながりが切れる分だけ位置エネルギーが増えるので、やはり $U$ は増えます（熱-1 ⑤の潜熱）。

内部エネルギーも、状態量です。系全体が動いているときの運動エネルギーや、系全体の高さによる位置エネルギー（力-4 の力学的エネルギー）は、$U$ にふくめず、別に数えます。ボールが転がっていても、中の空気の分子の熱運動のはげしさは変わらないからです。

%% ---------------------------------------------------------------
\bksubsection{3}{気体がする仕事：$p$-$V$ 図の面積}
\label{ph:熱-2:sec3}
\phboxhead{ピストンを押す力がする仕事}
\noindent\begin{minipage}[c]{0.46\linewidth}\centering
\begin{tikzpicture}[x=1mm, y=1mm]
  \qxcpiston{(0,0)}{48}{22}{26}
  \draw[bkInk, line width=0.5pt, dashed] (33,-1) -- (33,23);
  \draw[{Stealth[length=3pt]}-{Stealth[length=3pt]}, bkGray, line width=0.5pt] (26,-3) -- node[qxc lab, below] {$\dd x$} (33,-3);
  \draw[qxc arr, bkRed!80!black] (29.5,15) -- (40,15);
  \node[qxc lab, anchor=south west, text=bkRed!75!black] at (36,15.5) {$F=pA$};
  \node[qxc lab, text=bkBlue!80!black] at (12,11) {圧力 $p$};
  \node[qxc lab, anchor=south] at (28,23) {面積 $A$};
\end{tikzpicture}
\end{minipage}\hfill
\begin{minipage}[c]{0.51\linewidth}\raggedright
シリンダーの中の気体（圧力 $p$）は、面積 $A$ のピストンを、力 $F=pA$ で外へ押しています（圧力は1\,m$^2$あたりの力。力-12）。ピストンが外へ少し $\dd x$ 動くと、気体がピストンにする仕事は、力×動いた距離（力-4）です。
\end{minipage}

\par\vspace{4pt}
\[\delta W'=F\,\dd x=pA\,\dd x=p\,\dd V\]
$A\,\dd x$ は、気体の体積の増えた分 $\dd V$ です。$\delta$ は「小さな仕事」の印で、その意味は⑦で説明します。体積が $V_1$ から $V_2$ まで変わる間に気体がする仕事は、これをたした積分です（5-12）。
\pheq{W'=\int_{V_1}^{V_2}p\,\dd V}{J ＝ Pa × m$^3$ ＝ N/m$^2$ × m$^3$ ＝ N·m}
これが、数学 8-7⑩で「気体が外にする仕事」とよんだ量です。ふくらむ（$\dd V>0$）ときは正、縮む（$\dd V<0$）ときは負です。

\begin{seikaku}
$p$ を気体の圧力と考えてよいのは、ピストンが\textbf{ゆっくり}動いて、シリンダーの中の圧力がどこも同じに保たれるときです。このように、いつも「つり合いの状態」をたどる変化を\textbf{準静的な変化}といいます。\bkidx{じゅんせいてきなへんか}{準静的な変化}急に動かすと、ピストンの近くだけ圧力が違ったり、渦ができたりして、1つの $p$ では書けません。この章の $p$-$V$ 図の道は、すべて準静的な変化です。
\end{seikaku}

\phboxhead{圧力が一定なら $W'=p\,\Delta V$}
圧力が一定なら、$p$ を積分の外に出して
\pheq{W'=p\,(V_2-V_1)=p\,\Delta V}{J ＝ Pa × m$^3$}
\label{ph:熱-2:pdV}です。高校の物理の「気体が外にする仕事は、圧力×体積の変化」です。\qty{100}{kPa}で体積が\qty{1}{L}（$\qty{1e-3}{m^3}$）増えれば、$\num{100e3}\times\num{1e-3}=\qty{100}{J}$ です。

\begin{phflow}[気体がする仕事]
\phflowrow{$\delta W'=F\,\dd x=pA\,\dd x$}{準静的\\（中の圧力はどこも $p$）}{$W'=\displaystyle\int p\,\dd V$}
\phflowrow{$W'=\displaystyle\int_{V_1}^{V_2}p\,\dd V$}{圧力 $p$ が一定}{$W'=p\,\Delta V$}
\end{phflow}

\phboxhead{$p$-$V$ 図の面積と道すじ}
$W'=\int p\,\dd V$ は、$p$-$V$ 図で、道の下の面積です。8-7⑩では、同じ「1\,L・200\,kPa」から「3\,L・100\,kPa」への3つの道で、面積が 400\,J・200\,J・300\,J と違いました（図\ref{ph:熱-2:fig-3michi}）。\textbf{気体がする仕事は、はじめとおわりの状態だけでは決まらず、道すじで変わります}。

\begin{phbaisuru}{体積の変化 $\Delta V$}
圧力が一定なら、$W'=p\,\Delta V$ も2倍です。同じ大気圧の中で、風船を\qty{1}{L}ふくらませると、まわりの空気を押しのける仕事が約\qty{100}{J}、\qty{2}{L}なら約\qty{200}{J}要ります（例題3）。
\end{phbaisuru}

\phboxhead{表計算で、曲がった道の面積をたす}
1\,L・100\,kPa から 2\,L・50\,kPa まで、(A) $pV$ が一定（$pV=100$\,kPa·L。3-7⑥の注射器と同じ形）の曲がった道と、(B) まっすぐな道の、2つの道で気体がする仕事をくらべます。体積を\qty{0.1}{L}ずつに区切り、それぞれの区間を台形とみて面積をたします（8-7③・5-12②と同じ手順）。kPa×L＝J です。
\begin{bkcenter}
\bksheet{5}{21}{%
1 & $V$（L） & $p$（A） & 面積の和（A） & $p$（B） & 面積の和（B） \\
2 & 1.0 & 100.0 & 0.00 & 100.0 & 0.00 \\
3 & 1.1 & 90.9 & 9.55 & 95.0 & 9.75 \\
4 & 1.2 & 83.3 & 18.26 & 90.0 & 19.00 \\
5 & 1.3 & 76.9 & 26.27 & 85.0 & 27.75 \\
6 & 1.4 & 71.4 & 33.69 & 80.0 & 36.00 \\
7 & 1.5 & 66.7 & 40.59 & 75.0 & 43.75 \\
8 & 1.6 & 62.5 & 47.05 & 70.0 & 51.00 \\
9 & 1.7 & 58.8 & 53.12 & 65.0 & 57.75 \\
10 & 1.8 & 55.6 & 58.84 & 60.0 & 64.00 \\
11 & 1.9 & 52.6 & 64.25 & 55.0 & 69.75 \\
12 & 2.0 & 50.0 & 69.38 & 50.0 & 75.00 \\}
\bktabcaption{B3＝100/A3、C3＝C2＋(B2＋B3)/2*(A3－A2)、D3＝100－50*(A3－1)、E3＝E2＋(D2＋D3)/2*(A3－A2)。}\label{ph:熱-2:tab-sheet}
\end{bkcenter}
曲がった道（A）では約\qty{69.4}{J}、まっすぐな道（B）では\qty{75.0}{J}です。区切りを細かくすると、（A）は $100\ln2\approx\qty{69.3}{J}$ に近づきます（例題4）。はじめとおわりが同じでも、仕事は道で違います。
\begin{bkfigure}
\begin{tikzpicture}[x=1mm, y=1mm]
  \draw[bkInk, line width=0.7pt, -{Stealth[length=2mm]}] (0,0) -- (78,0) node[right, qxc lab] {$V$（L）};
  \draw[bkInk, line width=0.7pt, -{Stealth[length=2mm]}] (0,0) -- (0,50) node[above, qxc lab] {$p$（kPa）};
  \foreach \xx in {1,2,3} {\draw[bkInk, line width=0.5pt] ({\xx*22},0) -- ({\xx*22},-1.2); \node[qxc lab, anchor=north] at ({\xx*22},-1.5) {\xx};}
  \foreach \yy in {100,200} {\draw[bkInk, line width=0.5pt] (0,{\yy*0.2}) -- (-1.2,{\yy*0.2}); \node[qxc lab, anchor=east] at (-1.5,{\yy*0.2}) {\yy};}
  \fill[bkOrange!15] (22,0) -- (22,40) -- (66,40) -- (66,0) -- cycle;
  \draw[bkOrange!85!black, line width=1.4pt, -{Stealth[length=2.4mm]}] (22,40) -- (66,40) -- (66,20);
  \draw[bkBlue, line width=1.4pt, dash pattern=on 3pt off 1.5pt, -{Stealth[length=2.4mm]}] (22,40) -- (22,20) -- (66,20);
  \draw[bkPurple!80!black, line width=1.2pt, dash pattern=on 1pt off 1.2pt] (22,40) -- (66,20);
  \fill[bkInk] (22,40) circle (0.9) (66,20) circle (0.9);
  \node[qxc lab, anchor=south] at (44,41) {道1：400\,J};
  \node[qxc lab, anchor=north, fill=white] at (44,18.5) {道2：200\,J};
  \node[qxc lab, anchor=west, fill=white] at (46,33) {まっすぐ：300\,J};
  \node[qxc lab, anchor=south east] at (21,41) {はじめ};
  \node[qxc lab, anchor=west] at (67,22) {おわり};
\end{tikzpicture}
\bkfigcaption{8-7⑩と同じ3つの道（数学第8巻の図と同じ数）。気体がする仕事 $W'$ は道の下の面積で、道によって変わる（だいだいの面は、道1の下の面積）。}\label{ph:熱-2:fig-3michi}
\end{bkfigure}

%% ---------------------------------------------------------------
\bksubsection{4}{熱力学の第一法則：$\Delta U=Q+W$}
\label{ph:熱-2:sec4}
\phboxhead{エネルギーのバランス式}
系の内部エネルギーは、境界を通って入ってきた熱と仕事の分だけ増えます。熱-1 の熱のバランス式に、仕事を加えたものです。
\pheq{\Delta U=Q+W}{J ＝ J ＋ J}
\label{ph:熱-2:dai1}ここで、
\begin{bkcenter}
\fbox{\parbox{0.88\linewidth}{\raggedright
・$\Delta U$：系の内部エネルギーの変化（増えたら正）\\
・$Q$：系が\textbf{受けとった}熱（入ったら正、出ていったら負）\\
・$W$：系が\textbf{された}仕事（押し縮められたら正、ふくらんだら負）}}
\end{bkcenter}
です。これを\textbf{熱力学の第一法則}といいます。\bkidx{ねつりきがくのだいいちほうそく}{熱力学の第一法則}「エネルギーは、熱と仕事の形で出入りするが、勝手に生まれたり消えたりしない」という、エネルギー保存の法則を、熱までふくめて書いたものです。この本では、高校の物理の教科書と同じく、\textbf{$W$ は系が「される」仕事}と決めます。③の、気体が「する」仕事 $W'$ とは、符号が反対です。
\pheq{W=-W'=-\int_{V_1}^{V_2}p\,\dd V}{J ＝ Pa × m$^3$}
だから、第一法則は $\Delta U=Q-W'$ とも書けます。気体がふくらんで外へ仕事をすると（$W'>0$）、その分だけ内部エネルギーが減ります。

\phboxhead{矢印の図で数える}
符号を間違えないいちばん確かな方法は、\textbf{矢印の図}をかいて、入る量・出る量を、どれも正の数で書くことです。
\begin{bkfigure}
\begin{tikzpicture}[x=1mm, y=1mm]
  \filldraw[fill=bkBlue!12, draw=bkBlue!80!black, line width=0.9pt, rounded corners=2mm] (40,0) rectangle (80,26);
  \node[qxc lab, align=center, font=\sffamily\small] at (60,13) {系\\内部エネルギー $U$};
  \draw[qxc arr, bkOrange!85!black] (14,20) -- (39,20); \node[qxc lab, anchor=south, text=bkOrange!80!black] at (22,21) {入る熱 $Q_{\text{入}}$};
  \draw[qxc arr, bkOrange!85!black] (81,20) -- (106,20); \node[qxc lab, anchor=south, text=bkOrange!80!black] at (98,21) {出る熱 $Q_{\text{出}}$};
  \draw[qxc arr, bkPurple!80!black] (14,6) -- (39,6); \node[qxc lab, anchor=north, text=bkPurple!80!black] at (22,5) {される仕事 $W_{\text{入}}$};
  \draw[qxc arr, bkPurple!80!black] (81,6) -- (106,6); \node[qxc lab, anchor=north, text=bkPurple!80!black] at (98,5) {する仕事 $W_{\text{出}}$};
\end{tikzpicture}
\bkfigcaption{矢印の図。4つの量はどれも正の数。入る矢印はたし、出る矢印はひく。}\label{ph:熱-2:fig-yajirushi}
\end{bkfigure}
\pheq{\Delta U=Q_{\text{入}}-Q_{\text{出}}+W_{\text{入}}-W_{\text{出}}}{J}
$Q=Q_{\text{入}}-Q_{\text{出}}$、$W=W_{\text{入}}-W_{\text{出}}$ とすれば、$\Delta U=Q+W$ と同じです。この形なら、どちらの符号の決まりの本を読んでも、矢印を見て直せます。熱機関やサイクルの章（熱-5・熱-7・熱-8）では、この $W_{\text{出}}$・$W_{\text{入}}$ の書き方を使います。

\begin{phyomi}[$\Delta U=Q+W$]
「系のエネルギーが増えた分は、熱で入った分と、仕事で入った分の和」。お金の出入りと同じで、$U$ は財布の中身、$Q$ と $W$ は2つの口座からの入金（出金なら負の数）。財布の中身は、どの口座から入ったかを覚えていない。
\end{phyomi}

\phboxhead{小さな変化で書く}
状態が少しだけ変わるときは、
\pheq{\dd U=\delta Q+\delta W}{J}
と書きます。$\dd U$ は状態量 $U$ の小さな変化（全微分）ですが、$\delta Q$・$\delta W$ は「小さな熱」「小さな仕事」で、何かの量の変化ではありません。この違いは⑦で、ていねいに見ます。

\begin{seikaku}
\phboxhead{機械の本の書き方：$Q=\Delta U+W$}
エンジンやタービンを扱う機械工学の本では、機械が「する」仕事を正にしたいので、第一法則を
\[Q=\Delta U+W\qquad\text{（$W$ は系が\textbf{した}仕事）}\]
と書くことがよくあります。この本の $W'$ を $W$ と書いているだけで、中身は同じ法則です。
\begin{bkcenter}
\begin{tikzpicture}[x=1mm, y=1mm]
  \filldraw[fill=bkBlue!12, draw=bkBlue!80!black, line width=0.8pt, rounded corners=1.5mm] (14,0) rectangle (36,14);
  \node[qxc lab] at (25,7) {$\Delta U$};
  \draw[qxc arr, bkOrange!85!black] (0,10) -- (13,10); \node[qxc lab, anchor=south] at (5,11) {$Q$};
  \draw[qxc arr, bkPurple!80!black] (0,4) -- (13,4); \node[qxc lab, anchor=north east] at (13,3) {$W$（される）};
  \node[qxc lab, anchor=north, align=center] at (25,-6) {この本・高校\\$\Delta U=Q+W$};
  \filldraw[fill=bkBlue!12, draw=bkBlue!80!black, line width=0.8pt, rounded corners=1.5mm] (78,0) rectangle (100,14);
  \node[qxc lab] at (89,7) {$\Delta U$};
  \draw[qxc arr, bkOrange!85!black] (64,10) -- (77,10); \node[qxc lab, anchor=south] at (69,11) {$Q$};
  \draw[qxc arr, bkPurple!80!black] (101,4) -- (114,4); \node[qxc lab, anchor=north west] at (101,3) {$W$（する）};
  \node[qxc lab, anchor=north, align=center] at (89,-6) {機械の本\\$Q=\Delta U+W$};
\end{tikzpicture}
\end{bkcenter}
本を読むときは、まず「$W$ は、されるほうか、するほうか」を確かめましょう。この巻の記号カードにも書いておきます。
\end{seikaku}

\phboxhead{力学的エネルギーと内部エネルギー（力-4 の続き）}
力-4 ⑦では、摩擦などの保存力でない力がする仕事 $W_{\text{他}}$（力-4 では $W'$ と書きましたが、この巻では $W'$ を気体がする仕事に使います）の分だけ、力学的エネルギー $E=K+U_{\text{位}}$ が変わることを学びました（ここでは、位置エネルギーを内部エネルギー $U$ と区別するため $U_{\text{位}}$ と書きます）。
\[\Delta E=W_{\text{他}}\]
車がブレーキで止まるとき、摩擦力は車に負の仕事をして、力学的エネルギーを減らします。でも、ブレーキの円板とパッドは温まっています。ブレーキもふくめた全体を1つの系とみると、力学的エネルギーと内部エネルギーの和が、外から入る熱と仕事の分だけ変わります。
\pheq{\Delta E+\Delta U=Q+W}{J}
外との熱と仕事のやりとりを無視できるなら（$Q=0$、$W=0$）、$\Delta U=-\Delta E$ です。\textbf{減った力学的エネルギーが、そっくり内部エネルギーになって、温度を上げる}のです（例題6・例題7）。力-4 で「摩擦で失われた」と言ったエネルギーは、なくなったのではなく、ここにありました。

\phboxhead{電気の仕事も、仕事}
電気ストーブのヒーター線を系にふくめると、コンセントから入る電気のエネルギーは、境界を通って入る\textbf{仕事}（電気の仕事）です。ヒーター線の内部エネルギーが増えて熱くなり、そこから部屋へ\textbf{熱}が出ていきます。熱-1 ③で「ヒーターの熱 $Q=Pt$」と言ったのは、水だけを系にして、ヒーターを「まわり」とみたからです。\textbf{境界の引き方で、同じエネルギーが、熱とよばれたり仕事とよばれたりします}。だから、まず系を決めることが大切です。

\begin{motto}[3]
\phboxhead{1秒あたりの第一法則と、数学 6-4 の式}
時間とともに熱や仕事が出入りし続けるときは、第一法則を「1秒あたり」で書くと便利です。短い時間 $\Delta t$ の間の $\Delta U=Q+W$ の両辺を $\Delta t$ でわり、$\Delta t$ を0に近づけると（5-3）、
\[\frac{\dd U}{\dd t}=\dot Q+\dot W\]
です。$\dot Q$ は1秒あたりに入る熱〔W〕、$\dot W$ は1秒あたりにされる仕事（仕事率）〔W〕で、上の点は「1秒あたり」の印です（力-13 の $\dot m$ と同じ）。

電気ポット（ヒーター線と水を合わせて系にする）なら、$\dot W$ はヒーターの電力 $P$、$\dot Q$ は壁から逃げる熱 $-\dfrac{T-T_{\mathrm a}}{R_{\mathrm{th}}}$ です。水の温度が $\dd T$ 上がると $\dd U=C_{\mathrm{th}}\,\dd T$ なので、
\[C_{\mathrm{th}}\frac{\dd T}{\dd t}=P-\frac{T-T_{\mathrm a}}{R_{\mathrm{th}}}\]
となります。これは、数学 6-4⑥で「たまる熱＝入る熱－逃げる熱」から立てた微分方程式と、まったく同じ式です。6-4 の「ヒーターで温めながら、熱が逃げる」式は、第一法則を1秒あたりで書いたものだったのです。この「1秒あたりのバランス式」は、流れのある機械を扱う熱-4 で、もう一度使います。
\end{motto}

\begin{phflow}[第一法則から]
\phflowrow{$\dd U=\delta Q+\delta W$}{熱の出入りがない\\（断熱、$\delta Q=0$）}{$\Delta U=W$}
\phflowrow{$\dd U=\delta Q+\delta W$}{体積が変わらない\\（$\delta W=-p\,\dd V=0$）}{$\Delta U=Q$}
\phflowrow{$\Delta E+\Delta U=Q+W$}{外とのやりとりなし\\（$Q=0$、$W=0$）}{$\Delta U=-\Delta E$}
\end{phflow}

\begin{phbaisuru}{車の速さ}
止まるまでに減る運動エネルギー $\tfrac{1}{2}mv^2$ は4倍になり、ブレーキの円板の温度の上がりも、約4倍になります（逃げる熱を無視すると）。\qty{20}{m/s}で約\qty{22}{K}なら、\qty{40}{m/s}では約\qty{87}{K}です（例題6）。高速道路の長い下り坂でブレーキを使い続けると、ブレーキが熱くなりすぎて効かなくなることがある（フェード現象）のは、このためです。
\par\smallskip\noindent\begin{tikzpicture}
\begin{axis}[bkplot, width=96mm, height=42mm, xmin=0, xmax=44, ymin=0, ymax=100, xtick={10,20,30,40}, ytick={22,50,87}, xlabel={$v$（m/s）}, ylabel={$\Delta T$（K）}]
  \addplot[bkcurve, domain=0:42, samples=80] {600*x^2/(24*460)};
  \fill[bkRed] (axis cs:20,21.74) circle (1.8pt) (axis cs:40,86.96) circle (1.8pt);
  \node[qxc lab, anchor=north west, text=bkRed!75!black] at (axis cs:21,20) {速さ2倍で4倍};
  \node[qxc lab, anchor=north west, text=bkGray] at (axis cs:1,96) {車\qty{1200}{kg}、円板4枚で\qty{24}{kg}};
\end{axis}
\end{tikzpicture}
\end{phbaisuru}

%% ---------------------------------------------------------------
\bksubsection{5}{熱が出入りしないで押すと、温まる}
\label{ph:熱-2:sec5}
\phboxhead{断熱圧縮}
熱が出入りしない変化を\textbf{断熱変化}といいます。\bkidx{だんねつへんか}{断熱変化}断熱材で囲まなくても、\textbf{すばやく}変化させれば、熱が伝わる時間がないので、断熱に近くなります。熱が伝わるには時間がかかるからです（6-4 の時定数）。断熱変化では $Q=0$ なので、第一法則は
\pheq{\Delta U=W}{J ＝ J}
です。気体を押し縮めると、気体は仕事をされ（$W>0$）、その分だけ内部エネルギーが増えて、温度が上がります。これを\textbf{断熱圧縮}といいます。反対に、気体が断熱でふくらむと（\textbf{断熱膨張}）、外へした仕事の分だけ内部エネルギーが減り、温度が下がります。

\noindent\begin{minipage}[c]{0.36\linewidth}\centering
\begin{tikzpicture}[x=1mm, y=1mm]
  \foreach \px/\py/\ang in {6/5/30,14/12/150,9/18/250,20/6/80,24/16/200,16/22/320,4/12/10} {\fill[bkBlue!70] (\px,\py) circle (1.1); \draw[bkGray, -{Stealth[length=2.2pt]}, line width=0.5pt] (\px,\py) -- ++(\ang:3.5);}
  \draw[bkInk, line width=0.8pt] (32,0) -- (0,0) -- (0,26) -- (32,26);
  \filldraw[fill=bkGray!45, draw=bkInk, line width=0.7pt] (28,0.4) rectangle (31,25.6);
  \draw[qxc arr, bkPurple!80!black] (40,13) -- (32,13);
  \fill[bkBlue!70] (26,9) circle (1.1);
  \draw[bkRed, -{Stealth[length=2.6pt]}, line width=0.8pt] (26,9) -- ++(170:7);
  \node[qxc lab, anchor=north] at (16,-1) {動くピストンにはね返る};
\end{tikzpicture}
\end{minipage}\hfill
\begin{minipage}[c]{0.6\linewidth}\raggedright
分子の目で見ると、押しこまれてくるピストンにぶつかってはね返った分子は、ぶつかる前より速くなります。バットで打ち返したボールが、来たときより速く飛んでいくのと同じです。速い分子が増えるので、温度が上がります（くわしくは熱-14）。
\end{minipage}\par\vspace{4pt}

\phboxhead{\bkhatenaref の答え}
空気入れのピストンをいきおいよく押すと、筒の中の空気は、熱が逃げる間もなく押し縮められ（断熱圧縮）、手がした仕事の分だけ内部エネルギーが増えて、温度が上がります。空気がいちばん強く押し縮められるのは、ピストンが下まで来たときの、筒の先のほうです。タイヤの中の圧力より高くなるまで押し縮められた熱い空気が、何十回も筒の先にふれるので、筒の先が温まるのです。ピストンが擦れる摩擦の熱も少しはありますが、それなら、ピストンが動く筒の全体が同じように温まるはずです。

どれくらい熱くなるかを見積もってみましょう。20\,℃・1気圧の空気\qty{0.20}{L}（約\qty{0.24}{g}）を、すばやく $\tfrac{1}{3}$ の体積まで押し縮めると、手がする仕事は約\qty{28}{J}です（この値の求め方は熱-3 ⑥）。これが全部、空気の内部エネルギーになると、温度は約\qty{160}{K}も上がり、約180\,℃になる計算です（例題9）。でも、本物の空気入れでは、筒やホースに熱がすぐ逃げ、押すたびに新しい空気が入るので、筒の先は「温かい」くらいですみます。

\begin{phbaisuru}{押す仕事}
断熱なら $\Delta U=W$ なので、内部エネルギーの増えも2倍です。空気の内部エネルギーは温度にほぼ比例して増えるので（⑥）、温度の上がりも約2倍になります。
\end{phbaisuru}

\begin{phmijika}[雲ができる・ディーゼルエンジンが火花なしで燃える]
空気のかたまりが上昇すると、まわりの気圧が下がるので、ふくらみます。熱をほとんどやりとりしないので断熱膨張になり、温度が下がって、空気の中の水蒸気が水の粒になります。これが雲です（中学理科。熱-3 ⑥・熱-9）。反対に、ディーゼルエンジンは、空気をすばやく約 $\tfrac{1}{18}$ まで押し縮めて、約680\,℃にし、そこに燃料をふきこんで火花なしで燃やします（熱-3・熱-7）。スプレー缶を使い続けると缶が冷たくなるのも、中のガスがふくらむときに温度が下がるからです。
\end{phmijika}

%% ---------------------------------------------------------------
\bksubsection{6}{自由膨張：気体の $U$ は温度で決まる}
\label{ph:熱-2:sec6}
\noindent\begin{minipage}[c]{0.44\linewidth}\centering
\begin{tikzpicture}[x=1mm, y=1mm]
  \filldraw[fill=bkBlue!10, draw=bkInk, line width=0.7pt] (0,0) rectangle (52,24);
  \filldraw[fill=bkBlue!35, draw=bkInk, line width=0.7pt] (4,5) rectangle (20,19);
  \filldraw[fill=white, draw=bkInk, line width=0.7pt] (32,5) rectangle (48,19);
  \draw[bkInk, line width=1.4pt] (20,12) -- (32,12);
  \filldraw[fill=bkRed!50, draw=bkInk, line width=0.5pt] (26,12) circle (1.6);
  \node[qxc lab] at (12,12) {空気};
  \node[qxc lab] at (40,12) {真空};
  \node[qxc lab, anchor=south] at (26,14) {コック};
  \node[qxc lab, anchor=north] at (26,-1) {まわりの水の温度を測る};
\end{tikzpicture}
\end{minipage}\hfill
\begin{minipage}[c]{0.53\linewidth}\raggedright
ジュールは1845年、2つの銅の容器を管とコックでつなぎ、片方に約22気圧の空気を入れ、もう片方を真空にして、全体を水の中にしずめました。コックを開くと、空気は真空の容器へふき出して広がります。このとき、まわりの水の温度は、測れるほどには変わりませんでした。
\end{minipage}

\par\vspace{4pt}
第一法則で考えます。空気を系にすると、
\begin{itemize}[leftmargin=1.6\zw, itemsep=0pt, topsep=2pt]
\item 真空の中へ広がるとき、押しのける相手がいないので、外へ仕事をしません（$W=0$）。
\item 水の温度が変わらなかったので、熱の出入りもありません（$Q=0$）。
\end{itemize}
だから $\Delta U=0$ です。このように、真空へ断熱で広がる変化を\textbf{断熱自由膨張}といいます。\bkidx{だんねつじゆうぼうちょう}{断熱自由膨張}
\pheq{\Delta U=Q+W=0+0=0}{J}
\label{ph:熱-2:jiyuu}自由膨張のあと、空気の体積は約2倍、圧力は約半分になりましたが、温度は変わりませんでした。\textbf{体積や圧力が変わっても、温度が同じなら内部エネルギーは同じ}だったのです。このことから、気体（とくに、薄い気体）の内部エネルギーは、温度だけで決まることがわかります。この性質をきちんともつ気体を\textbf{理想気体}といいます（熱-3）。高校の物理の「断熱自由膨張では $\Delta U=0$ で、理想気体の温度は変わらない」は、このことです。

\begin{seikaku}
もっと精密な実験では、ふつうの気体は、自由膨張でわずかに温度が下がります。分子どうしが弱く引き合っているので、分子が離れると位置エネルギーが増え、その分だけ運動エネルギーが減るからです。このずれは、圧力が低いほど小さくなります。ずれのくわしい扱いは熱-10 で学びます。また、自由膨張は、ひとりでに逆もどりしない（空気が勝手に元の容器にもどることはない）変化です。そのわけは、熱-6 ④でエントロピーを使って説明します。
\end{seikaku}

%% ---------------------------------------------------------------
\bksubsection{7}{状態量と、道すじで変わる量}
\label{ph:熱-2:sec7}
\phboxhead{$U$ は道によらない、$Q$ と $W$ は道による}
⑥までで、次のことがわかりました。
\begin{itemize}[leftmargin=1.6\zw, itemsep=0pt, topsep=2pt]
\item 内部エネルギー $U$ は状態量なので、はじめとおわりの状態が同じなら、どんな道を通っても $\Delta U$ は同じ。
\item 仕事 $W$ は、$p$-$V$ 図の道で変わる（③）。
\item だから、$Q=\Delta U-W$ も道で変わる（例題5）。
\end{itemize}
数学 8-3 の専門書は、このことを「$\dd U$ は完全微分であるが、熱 $\delta Q$ と仕事 $\delta W$ は完全微分ではない」と書いていました。

\phboxhead{全微分かどうかの見分け方}
2つの変数 $x$、$y$ の小さな変化で書いた量 $A\,\dd x+B\,\dd y$（$A$、$B$ は $x$、$y$ の関数）が、ある関数 $f(x,\ y)$ の全微分 $\dd f$ になっているとします。全微分（8-3⑤）は
\[\dd f=\frac{\partial f}{\partial x}\,\dd x+\frac{\partial f}{\partial y}\,\dd y\]
なので、$A=\partial f/\partial x$、$B=\partial f/\partial y$ です。2回の偏微分は順番を変えても同じ（8-2⑧）なので、
\[\frac{\partial A}{\partial y}=\frac{\partial B}{\partial x}\qquad\text{（全微分なら必ず成り立つ）}\]
でなければなりません。反対に、穴のない広がりの中でこれが成り立てば、$A\,\dd x+B\,\dd y$ は全微分です（力-4 の保存力の見分け方 $\mathrm{rot}\,\vec F=\vec 0$ と同じ。8-10）。全微分になっている量を\textbf{完全微分}といいます。\bkidx{かんぜんびぶん}{完全微分}

\phboxhead{気体の仕事に当てはめる}
状態を $(V,\ p)$ の2つの変数で表します。気体がする小さな仕事は $\delta W'=p\,\dd V+0\,\dd p$ です。$A=p$、$B=0$ として、
\[\frac{\partial A}{\partial p}=\frac{\partial p}{\partial p}=1,\qquad \frac{\partial B}{\partial V}=0\]
で、等しくありません。だから $p\,\dd V$ は完全微分ではなく、仕事は道で変わります。8-7⑩の3つの道で仕事が違ったのは、このためです。一方、$p\,\dd V+V\,\dd p$ なら $A=p$、$B=V$ で、$1=1$ です。これは $pV$ の全微分 $\dd(pV)$ で、道によらず「おわりの $pV$ －はじめの $pV$」で決まります。

\begin{phyomi}[$\dd U$ と $\delta Q$]
「$\dd$ は『状態量の小さな変化』、$\delta$ は『小さな受けわたしの量』」。山の高さの変化（$\dd h$）は歩く道によらないが、歩いた道のり（$\delta$）は道で変わるのと同じ（8-3）。
\end{phyomi}

\phboxhead{体積一定なら、熱は $U$ の変化と同じ}
体積を変えないで温めると、$\delta W=0$ なので $\dd U=\delta Q$ です。このときだけは、受けとった熱が状態量 $U$ の変化と等しく、道によらなくなります。体積一定で温度を\qty{1}{K}上げるのに要る熱（定積熱容量）は、8-3 の専門書の書き方で
\[C_V=\Bigl(\frac{\partial U}{\partial T}\Bigr)_V\]
と書けます（右下の $V$ は「体積を止めて」）。圧力一定で温めると、気体はふくらんで外へ仕事をするので、同じ\qty{1}{K}上げるのに、その分だけ多くの熱が要ります。この2つの比熱（定積比熱と定圧比熱）の違いは、熱-3 ④で求めます。

%% ---------------------------------------------------------------
\Needspace{14\baselineskip}
\bksubsection{8}{1周してもどる：サイクル}
\label{ph:熱-2:sec8}
\noindent\begin{minipage}[c]{0.44\linewidth}\centering
\begin{tikzpicture}[x=0.76mm, y=1mm]
  \draw[bkInk, line width=0.7pt, -{Stealth[length=2mm]}] (0,0) -- (56,0) node[right, qxc lab] {$V$（L）};
  \draw[bkInk, line width=0.7pt, -{Stealth[length=2mm]}] (0,0) -- (0,40) node[above, qxc lab] {$p$（kPa）};
  \foreach \xx in {1,2,3} {\draw[bkInk, line width=0.5pt] ({\xx*15},0) -- ({\xx*15},-1.2); \node[qxc lab, anchor=north] at ({\xx*15},-1.5) {\xx};}
  \foreach \yy in {100,200} {\draw[bkInk, line width=0.5pt] (0,{\yy*0.15}) -- (-1.2,{\yy*0.15}); \node[qxc lab, anchor=east] at (-1.5,{\yy*0.15}) {\yy};}
  \fill[bkOrange!20] (15,15) rectangle (45,30);
  \draw[bkOrange!85!black, line width=1.3pt, -{Stealth[length=2.4mm]}] (15,30) -- (45,30);
  \draw[bkOrange!85!black, line width=1.3pt, -{Stealth[length=2.4mm]}] (45,30) -- (45,15);
  \draw[bkOrange!85!black, line width=1.3pt, -{Stealth[length=2.4mm]}] (45,15) -- (15,15);
  \draw[bkOrange!85!black, line width=1.3pt, -{Stealth[length=2.4mm]}] (15,15) -- (15,30);
  \node[qxc lab, text=bkOrange!75!black] at (30,22.5) {200\,J};
  \node[qxc lab, anchor=south] at (15,31) {A}; \node[qxc lab, anchor=south] at (45,31) {B};
  \node[qxc lab, anchor=north west] at (45.5,14.5) {C}; \node[qxc lab, anchor=north east] at (14.5,14.5) {D};
\end{tikzpicture}
\end{minipage}\hfill
\begin{minipage}[c]{0.53\linewidth}\raggedright
気体の状態が、いくつかの変化をへて、はじめの状態にもどることを\textbf{サイクル}といいます。\bkidx{さいくる}{サイクル}左の図は、A→B（200\,kPa でふくらむ）→C（体積一定で圧力を下げる）→D（100\,kPa で縮む）→A（体積一定で圧力を上げる）と、時計まわりに1周するサイクルです。
\end{minipage}

\par\vspace{4pt}
1周すると、状態は元にもどるので、状態量 $U$ も元の値です。1周の積分（8-7 の $\oint$）で書くと
\[\oint\dd U=0\]
です。1周の間に出入りした量を、矢印の図（④）でたし合わせると
\[0=Q_{\text{入}}-Q_{\text{出}}+W_{\text{入}}-W_{\text{出}}\]
です。つまり
\pheq{W_{\text{出}}-W_{\text{入}}=Q_{\text{入}}-Q_{\text{出}}}{J ＝ J}
です。\textbf{1周で気体が外にした正味の仕事は、1周で受けとった正味の熱に等しい}。熱を受けとって、それを仕事に変え続ける機械（熱機関）は、このしくみで動いています（熱-5）。

正味の仕事は、$p$-$V$ 図の上で、サイクルが囲む面積です。A→B でふくらむときに外へする仕事（$200\times2=\qty{400}{J}$）から、C→D で縮むときにされる仕事（$100\times2=\qty{200}{J}$）をひいた\qty{200}{J}で、ちょうど長方形の面積です（B→C と D→A は体積が変わらないので、仕事は0。例題12）。時計まわりのサイクルでは $W_{\text{出}}-W_{\text{入}}>0$（外へ仕事をする）、反時計まわりでは $W_{\text{出}}-W_{\text{入}}<0$（外から仕事をされる）です。8-7 の終わりで「エンジンの気体が1回にする仕事は、閉じた道が囲む面積」と書いたのは、この正味の仕事のことです（熱-6 ⑧で、温度とエントロピーの図でも同じことを見ます）。

\begin{phmada}[小さな小さな機械では、熱と仕事をどう分ける？]
この章の熱と仕事は、たくさんの分子をまとめて考えたときの量です。ところが、体の中で荷物を運ぶモータータンパク質（1個の大きな分子でできた、ナノメートルの大きさの機械）のような小さな系では、まわりの分子がぶつかるゆらぎ（熱ゆらぎ）がとても大きく、仕事と熱の量が1回ごとにばらつきます。こうした小さな系で熱と仕事をどう決めて数えるかは、1990年代の終わりごろから「確率熱力学」として研究が進んでいます（日本の関本謙らの研究が先がけの1つ）。熱ゆらぎから仕事を取り出す実験もおこなわれていますが、生き物の中の機械が、どれくらい効率よく熱ゆらぎを使っているのかは、まだよくわかっていません。
\end{phmada}

\begin{phsaki}{化学反応の熱}
燃料が燃える熱（熱-1 ⑧）や、物がとけるときの熱は、化学反応で内部エネルギーが変わることから生まれます。反応の熱を、反応の道すじによらない状態量の差として計算する方法（ヘスの法則）は、化学の熱力学で学びます。圧力一定の反応の熱は、内部エネルギーではなく「エンタルピー」の差になります。エンタルピーは、熱-4 で流れる気体のために作ります。
\end{phsaki}

\begin{senmonsho}[機械工学の熱力学の本の書き方]
\begin{bkhonbun}
閉じた系が状態1から状態2へ変化するとき，系に加えられた熱量を $Q_{12}$，系が外部になした仕事を $W_{12}$ とすると，$Q_{12}=U_2-U_1+W_{12}$ が成り立つ．準静的過程では $W_{12}=\int_1^2p\,dV$ であり，$\delta Q=dU+p\,dV$ と書ける．
\end{bkhonbun}
\bkyaku{「閉じた系」は、物が出入りしない系（②）。「なした仕事」は「した仕事」で、この本の $W'$ にあたる（$W_{12}=W'=-W$）。だから、この文の式は、この本の $\Delta U=Q+W$ と同じ。$Q_{12}$ の添字12は「状態1から2までの間に」の意味。$\delta Q=dU+p\,dV$ は、この本の $\dd U=\delta Q+\delta W$ に $\delta W=-p\,\dd V$ を入れて並べかえたもの（④・③）。}
\end{senmonsho}

\begin{taisetsu}
\def\tsl#1{\makebox[6.8\zw][l]{\sffamily\bfseries\small\color{\bkTT}#1}}
\tsl{熱と仕事}どちらもエネルギーの移り方。\qty{1}{cal}＝\qty{4.184}{J}\\[3pt]
\tsl{気体の仕事}$W'=\displaystyle\int p\,\dd V$（$p$-$V$ 図の面積）。$p$ 一定で $W'=p\,\Delta V$\\[3pt]
\tsl{第一法則}$\Delta U=Q+W$（$W$ は「される」仕事。$W=-W'$）\\[3pt]
\tsl{矢印の図}$\Delta U=Q_{\text{入}}-Q_{\text{出}}+W_{\text{入}}-W_{\text{出}}$\\[3pt]
\tsl{断熱など}断熱 $\Delta U=W$。体積一定 $\Delta U=Q$。自由膨張 $\Delta U=0$\\[3pt]
\tsl{状態量}$U$ は道によらない（完全微分 $\dd U$）。$Q$・$W$ は道による（$\delta Q$・$\delta W$）\\[3pt]
\tsl{サイクル}$\oint\dd U=0$。正味の仕事 $W_{\text{出}}-W_{\text{入}}$＝正味の熱＝囲む面積
\end{taisetsu}

\begin{kotoba}
\small
\bkword{熱の仕事当量}{熱と仕事の換算の比。\qty{1}{cal}＝\qty{4.184}{J}。}
\bkword{系・まわり・境界}{考える物のまとまり、そのほか、2つを分ける面。}
\bkword{状態量}{いまの状態だけで決まる量（$p$・$V$・$T$・$U$）。}
\bkword{内部エネルギー}{系の中の分子の運動エネルギーと、分子どうしの位置エネルギーの和。}
\bkword{準静的な変化}{いつもつり合いの状態をたどる、ゆっくりした変化。}
\bkword{熱力学の第一法則}{$\Delta U=Q+W$。熱もふくめたエネルギー保存の法則。}
\bkword{断熱変化・断熱自由膨張}{熱の出入りがない変化。真空へ広がる断熱の変化。}
\bkword{完全微分}{ある関数の全微分になっている量。道によらずに積分できる。}
\bkword{サイクル}{いくつかの変化をへて、はじめの状態にもどること。}
\end{kotoba}

\bkkigocards{\bkkigocard{$U$}{ユー}{内部エネルギー〔J〕。状態量。この章では、力-4 の位置エネルギーを $U_{\text{位}}$ と書く}{internal energy}{熱-2}}{\bkkigocard{$W$}{ダブリュー}{系が\textbf{される}仕事〔J〕。$\Delta U=Q+W$。機械の本では「した仕事」に使う}{work (done on)}{熱-2}}
\bkkigocards{\bkkigocard{$W'$}{ダブリュー・ダッシュ}{系が\textbf{する}仕事〔J〕。$W'=-W$。サイクルの章では $W_{\text{出}}$・$W_{\text{入}}$ とも書く}{work (done by)}{熱-2}}{\bkkigocard{$\delta Q$}{デルタ・キュー}{小さな熱。状態量の変化ではないので $\dd$ でなく $\delta$。$\delta W$ も同じ}{small heat}{熱-2}}

%% ======================== 例題 ========================
\Needspace{10\baselineskip}\bksection{reidai}{例題でたしかめよう}
\noindent 水の比熱は 4.2\,kJ/(kg·K)、$g=\qty{9.8}{m/s^2}$ とします。答えは2けたで書きます。符号は、この章の決まり（$\Delta U=Q+W$、$W$ は系が\textbf{される}仕事、$W'=-W$ は系が\textbf{する}仕事）で書きます。

\label{ph:熱-2:ex1}
\begin{phreidai}[ジュールの実験：水はどれだけ温まる？]
質量\qty{10}{kg}の重りが\qty{2.0}{m}落ちて、羽根車で水\qty{1.0}{kg}をかき回します。これを20回繰り返すと、水の温度は何 K 上がりますか。
\phbamen
重りの位置エネルギーが、羽根車の仕事になり、水の内部エネルギーを増やします（①の実験）。
\phmushi
重りが下に着くときの運動エネルギー、ひもと滑車の摩擦、羽根車と容器が受けとる熱、まわりへ逃げる熱を無視します。
\phzu
系は水。重り $m=\qty{10}{kg}$、高さ $h=\qty{2.0}{m}$、回数20、水 $M=\qty{1.0}{kg}$、$c=4.2$\,kJ/(kg·K)。
\phriyu
水は熱を受けとらず（$Q=0$）、羽根車から仕事をされるだけなので、第一法則（④）は $\Delta U=W$ です。水の内部エネルギーの増えは、温度の上がりで $Mc\,\Delta T$ と書けます（熱-1 ③）。
\phtateru
\[W=20\times mgh,\qquad Mc\,\Delta T=W\]
\phtoku
\[W=20\times10\times9.8\times2.0=\qty{3920}{J},\qquad \Delta T=\frac{3920}{1.0\times4200}\approx\qty{0.93}{K}\]
\phtashikame
大きさ：重りを40\,m分も落としても、1\,K に届きません。熱にすると、仕事はとても小さく見えます。ジュールの実験が難しかったのは、このためです。\\
単位：J ÷ (kg × J/(kg·K)) ＝ K。
\kotae 約\qty{0.93}{K}
\end{phreidai}
\begin{jibun}
質量\qty{25}{kg}の重りが\qty{1.5}{m}落ちるのを30回繰り返して、水\qty{2.0}{kg}をかき回します。水の温度の上がりを求めましょう。
\bkjans{$W=30\times25\times9.8\times1.5=\qty{11025}{J}$。$11025\div(2.0\times4200)$ より、約\qty{1.3}{K}}
\end{jibun}

\label{ph:熱-2:ex2}
\begin{phreidai}[食事のエネルギーと、体が出す熱]
大人が1日に食べ物からとるエネルギーを\qty{2000}{kcal}とします。これは何 J ですか。このエネルギーが、1日かけて一定の割合で体から出ていくとすると、何 W にあたりますか（\qty{1}{cal}＝\qty{4.184}{J}）。
\phbamen
食べ物の化学エネルギーが、体の中で、筋肉の仕事と熱になります。体温が毎日同じなら、入った分だけ出ていきます。
\phmushi
体重の変化（たくわえられる分）を無視します。体がする仕事も、最後は熱になってまわりへ出るとします。
\phzu
1日 $t=24\times3600=\qty{86400}{s}$。エネルギー $E=\qty{2000}{kcal}$。
\phriyu
cal から J への換算（熱の仕事当量、①）と、「1秒あたりのエネルギー＝仕事率」（力-4 ⑧）を使います。体を系にすると、1日で $\Delta U\approx0$ なので、入る分＝出る分です。
\phtateru
\[E=2000\times10^3\times4.184\ \text{J},\qquad P=\frac{E}{t}\]
\phtoku
\[E\approx\qty{8.37e6}{J}=\qty{8.4}{MJ},\qquad P=\frac{\num{8.37e6}}{86400}\approx\qty{97}{W}\]
\phtashikame
大きさ：人の体は、ふだん約\qty{100}{W}の電球ほどの熱を出しています。満員の部屋が暑くなるのは、このためです（冷房の設計では、人1人あたりの熱を数えます）。\\
単位：kcal は「k（1000）× cal」なので、1000 をかけ忘れないようにします。
\kotae 約\qty{8.4}{MJ}、約\qty{97}{W}
\end{phreidai}
\begin{jibun}
1日\qty{2400}{kcal}をとる人の、平均の仕事率（体が出す熱）を求めましょう。
\bkjans{$2400\times10^3\times4.184\div86400$ より、約\qty{1.2e2}{W}（116）}
\end{jibun}

\label{ph:熱-2:ex3}
\begin{phreidai}[1気圧でふくらむ気体の仕事]
ピストンのついたシリンダーの中の気体を温めたら、圧力\qty{101.3}{kPa}（1気圧）のまま、体積が\qty{2.0}{L}から\qty{3.0}{L}にふくらみました。気体がした仕事 $W'$ と、気体がされた仕事 $W$ を求めましょう。
\phbamen
ピストンの外は大気で、気体は大気を押しのけながらふくらみます。
\phmushi
ピストンの摩擦と質量を無視します。ゆっくりふくらむ（準静的）とします。
\phzu
$p=\qty{101.3e3}{Pa}$、$\Delta V=\qty{1.0}{L}=\qty{1.0e-3}{m^3}$。
\phriyu
圧力が一定なので、③の $W'=p\,\Delta V$ が使えます。$W$ は $W'$ の符号を変えたもの（④）です。
\phtateru
\[W'=p\,\Delta V,\qquad W=-W'\]
\phtoku
\[W'=\num{101.3e3}\times\num{1.0e-3}\approx\qty{1.0e2}{J},\qquad W\approx-\qty{1.0e2}{J}\]
\phtashikame
単位：Pa × m$^3$ ＝ N/m$^2$ × m$^3$ ＝ N·m ＝ J。\\
大きさ：\qty{100}{J}は、\qty{1}{kg}の物を約\qty{10}{m}持ち上げる仕事です。\\
符号：ふくらむので気体は外へ仕事をし、$W'>0$、$W<0$。
\kotae $W'\approx\qty{1.0e2}{J}$（\qty{101}{J}）、$W\approx-\qty{1.0e2}{J}$
\end{phreidai}
\begin{jibun}
圧力\qty{200}{kPa}のまま、気体の体積が\qty{0.50}{L}から\qty{1.5}{L}にふくらみました。$W'$ と $W$ を求めましょう。
\bkjans{$W'=\num{200e3}\times\num{1.0e-3}=\qty{200}{J}$、$W=-\qty{200}{J}$}
\end{jibun}

\label{ph:熱-2:ex4}
\begin{phreidai}[$pV$ が一定の道で、気体がする仕事]
はじめ\qty{100}{kPa}・\qty{1.0}{L}の気体が、$pV$ を一定に保ったまま（3-7⑥の注射器と同じ形）、\qty{2.0}{L}までふくらみました。気体がした仕事を求めましょう。
\phbamen
圧力が一定ではないので、$p\,\Delta V$ は使えません。$p$-$V$ 図の曲線の下の面積を求めます。
\phmushi
準静的な変化とします。$pV=100$\,kPa·L はちょうど一定とします。
\phzu
$p=\dfrac{100}{V}$（$p$ は kPa、$V$ は L）。$V$ は 1.0 から 2.0 まで。kPa × L ＝ J。
\phriyu
③の $W'=\int p\,\dd V$ に、$p$ を $V$ の式で入れます。$\int\frac{1}{x}\,\dd x=\ln x$（5-14⑨）を使います。
\phtateru
\[W'=\int_{1.0}^{2.0}\frac{100}{V}\,\dd V\]
\phtoku
\[W'=100\,\bigl[\ln V\bigr]_{1.0}^{2.0}=100\,(\ln2.0-\ln1.0)=100\times0.693\approx\qty{69}{J}\]
\phtashikame
表計算：③の表（A）で、\qty{0.1}{L}ずつの台形をたした値は\qty{69.38}{J}で、ほぼ同じです。\\
くらべる：同じはじめとおわりを、まっすぐな道で結ぶと\qty{75}{J}（③の表（B））。曲がった道のほうが、下の面積が小さい。\\
2倍にすると：体積を4倍（\qty{4.0}{L}）までふくらませると、$100\ln4=2\times100\ln2$ で、仕事は2倍の約\qty{139}{J}です。
\kotae 約\qty{69}{J}
\end{phreidai}
\begin{jibun}
はじめ\qty{200}{kPa}・\qty{1.0}{L}の気体が、$pV$ を一定に保ったまま\qty{3.0}{L}までふくらみました。気体がした仕事を求めましょう（$\ln3=1.099$）。
\bkjans{$200\ln3=200\times1.099$ より、約\qty{2.2e2}{J}}
\end{jibun}

\label{ph:熱-2:ex5}
\begin{phreidai}[8-7 の3つの道：受けとる熱は道で変わる]
ある気体が、1\,L・200\,kPa の状態から 3\,L・100\,kPa の状態へ、8-7⑩の3つの道（図\ref{ph:熱-2:fig-3michi}）で変わります。この気体では、2つの状態の内部エネルギーの差が $\Delta U=\qty{250}{J}$ とわかっています（この値の出し方は熱-3 ④）。それぞれの道で、気体が受けとる熱 $Q$ を求めましょう。
\phbamen
はじめとおわりは同じで、道だけが違います。
\phmushi
どの道も準静的な変化とします。
\phzu
気体がする仕事は、道1 $W'=\qty{400}{J}$、道2 $W'=\qty{200}{J}$、まっすぐ $W'=\qty{300}{J}$（8-7⑩）。どれも $\Delta U=\qty{250}{J}$。
\phriyu
$U$ は状態量なので、どの道でも $\Delta U$ は同じです（⑦）。第一法則 $\Delta U=Q+W$ を $Q$ について解き、$W=-W'$ を入れます（④）。
\phtateru
\[Q=\Delta U-W=\Delta U+W'\]
\phtoku
\[\text{道1：}Q=250+400=\qty{650}{J}\qquad\text{道2：}Q=250+200=\qty{450}{J}\]
\[\text{まっすぐ：}Q=250+300=\qty{550}{J}\]
\phtashikame
矢印の図で：どの道も、入る熱の一部（\qty{250}{J}）が内部エネルギーにたまり、残りが外への仕事として出ていきます。\\
道による：$\Delta U$ は同じでも、$Q$ と $W'$ は道で変わります。ただし、$Q-W'$ はどの道でも\qty{250}{J}です。
\kotae 道1 \qty{650}{J}、道2 \qty{450}{J}、まっすぐ \qty{550}{J}
\end{phreidai}
\begin{jibun}
同じ気体を、3\,L・100\,kPa から、圧力\qty{100}{kPa}のまま\qty{1}{L}まで縮め、そのあと体積を変えずに\qty{200}{kPa}まで上げました（道2の逆）。$W$ と $Q$ を求めましょう（$\Delta U=-\qty{250}{J}$）。
\bkjans{$W=+\qty{200}{J}$（縮むので、された仕事は正）。$Q=\Delta U-W=-250-200=-\qty{450}{J}$（\qty{450}{J}の熱を出す）}
\end{jibun}

\label{ph:熱-2:ex6}
\begin{phreidai}[ブレーキの円板は、何 K 温まる？（力-4 の続き）]
質量\qty{1200}{kg}の車が、\qty{20}{m/s}（時速\qty{72}{km}）から、ブレーキで止まりました。止まる間に、ブレーキの円板4枚（合わせて\qty{24}{kg}、比熱 0.46\,kJ/(kg·K)）の温度は何 K 上がりますか。
\phbamen
車の運動エネルギーが、ブレーキの摩擦で、円板とパッドの内部エネルギーになります（④の「力学的エネルギーと内部エネルギー」）。
\phmushi
止まる間に、円板から空気へ逃げる熱を無視します。パッドとタイヤが受けとる分も無視し、全部が円板に入るとします。道は水平です。
\phzu
系は車全体（円板をふくむ）。$m=\qty{1200}{kg}$、$v=\qty{20}{m/s}$、円板 $M=\qty{24}{kg}$、$c=460$\,J/(kg·K)。
\phriyu
外との熱と仕事のやりとりを無視するので、④の $\Delta U=-\Delta E$ です。減る力学的エネルギーは運動エネルギー $\tfrac{1}{2}mv^2$（力-4 ⑤）、増える内部エネルギーは $Mc\,\Delta T$ です。
\phtateru
\[Mc\,\Delta T=\frac{1}{2}mv^2\]
\phtoku
\[\frac{1}{2}mv^2=\frac{1}{2}\times1200\times20^2=\qty{240000}{J}\]
\[\Delta T=\frac{\num{240000}}{24\times460}\approx\qty{22}{K}\]
\phtashikame
大きさ：1回の停止で\qty{20}{K}ほど上がります。何度も繰り返したり、長い下り坂でブレーキを使い続けたりすると、円板は数百℃にもなります。\\
2倍にすると：速さが2倍なら、温度の上がりは4倍（④の「2倍にすると」）。
\kotae 約\qty{22}{K}
\end{phreidai}
\begin{jibun}
質量\qty{1500}{kg}の車が\qty{25}{m/s}から止まります。ブレーキの円板4枚（合わせて\qty{32}{kg}、比熱 0.46\,kJ/(kg·K)）の温度の上がりを求めましょう（逃げる熱は無視）。
\bkjans{$\tfrac{1}{2}\times1500\times25^2=\qty{468750}{J}$。$468750\div(32\times460)$ より、約\qty{32}{K}}
\end{jibun}

\label{ph:熱-2:ex7}
\begin{phreidai}[自転車で坂を下る：ブレーキで熱になる量]
人と自転車を合わせて\qty{70}{kg}が、高さ\qty{100}{m}の坂を、ブレーキをかけながら一定の速さで下りました。ブレーキで熱（内部エネルギー）になったエネルギーを求めましょう。
\phbamen
速さが変わらないので、運動エネルギーは変わりません。減るのは位置エネルギーだけです。
\phmushi
空気の抵抗とタイヤの転がる抵抗を無視します（あれば、その分ブレーキの熱が減る）。
\phzu
$m=\qty{70}{kg}$、下りた高さ $h=\qty{100}{m}$、$\Delta K=0$。
\phriyu
④の $\Delta U=-\Delta E$ で、$\Delta E=\Delta K+\Delta U_{\text{位}}=0-mgh$ です。
\phtateru
\[\Delta U=-\Delta E=mgh\]
\phtoku
\[\Delta U=70\times9.8\times100=\qty{68600}{J}\approx\qty{69}{kJ}\]
\phtashikame
大きさ：もし、この熱が全部、アルミニウムの車輪のリム2本（合わせて\qty{1.0}{kg}、比熱 0.90\,kJ/(kg·K)）にたまったら、$68.6\div0.90\approx\qty{76}{K}$ も上がります。本物は、走る風でたえず冷やされるので、そこまでは上がりません。\\
力-4 とくらべる：力-4 で「摩擦の負の仕事の分、力学的エネルギーが減る」と言った量が、そのまま内部エネルギーになっています。
\kotae 約\qty{69}{kJ}
\end{phreidai}
\begin{jibun}
人と自転車で\qty{80}{kg}が、高さ\qty{50}{m}の坂を、ブレーキをかけながら一定の速さで下りました。ブレーキで熱になったエネルギーを求めましょう。
\bkjans{$80\times9.8\times50=\qty{39200}{J}$ より、約\qty{39}{kJ}}
\end{jibun}

\label{ph:熱-2:ex8}
\begin{phreidai}[ミキサーでかき混ぜると、水は温まる]
断熱した入れ物の中の水\qty{1.0}{kg}を、\qty{300}{W}のミキサーで\qty{60}{s}かき混ぜました。モーターがした仕事が全部、羽根から水に入るとして、水の温度の上がりを求めましょう。
\phbamen
羽根が水をかき回す仕事が、水の内部エネルギーになります（ジュールの実験と同じしくみ）。
\phmushi
モーターの中でむだになる分は、ないとします（全部が羽根の仕事になる）。入れ物と羽根の熱容量を無視します。
\phzu
系は水。$P=\qty{300}{W}$、$t=\qty{60}{s}$、$M=\qty{1.0}{kg}$。
\phriyu
断熱なので $Q=0$、第一法則は $\Delta U=W$（⑤）。仕事は「仕事率×時間」（力-4 ⑧）。
\phtateru
\[W=Pt,\qquad Mc\,\Delta T=W\]
\phtoku
\[W=300\times60=\qty{18000}{J},\qquad \Delta T=\frac{18000}{1.0\times4200}\approx\qty{4.3}{K}\]
\phtashikame
くらべる：ジュールの実験（例題1）より、ずっと大きく温まります。モーターは、重りよりずっと大きな仕事率を出せるからです。\\
2倍にすると：時間を2倍にすると、温度の上がりも2倍（約\qty{8.6}{K}）。
\kotae 約\qty{4.3}{K}
\end{phreidai}
\begin{jibun}
断熱した入れ物の中の水\qty{2.0}{kg}を、\qty{500}{W}のミキサーで\qty{120}{s}かき混ぜます。水の温度の上がりを求めましょう。
\bkjans{$W=500\times120=\qty{60000}{J}$。$60000\div(2.0\times4200)$ より、約\qty{7.1}{K}}
\end{jibun}

\label{ph:熱-2:ex9}
\begin{phreidai}[空気入れの空気は、何℃になる？（\bkhatenaref）]
20\,℃・1気圧の空気\qty{0.20}{L}を、空気入れですばやく $\tfrac{1}{3}$ の体積まで押し縮めたら、手がした仕事は約\qty{28}{J}でした。熱が出入りしないとすると、空気の温度は何 K 上がりますか。空気の密度は\qty{1.20}{kg/m^3}、体積を変えないで\qty{1}{K}上げる熱（定積比熱）は $c_v=718$\,J/(kg·K) とします（熱-3 ④で使う値）。
\phbamen
すばやく押すので、熱が逃げる時間がありません（⑤の断熱圧縮）。
\phmushi
筒・ホースへ逃げる熱、ピストンの摩擦、空気がタイヤへ流れ出すことを無視します（弁が開く前の1回の押し縮め）。
\phzu
系は筒の中の空気。$V=\qty{0.20}{L}=\qty{2.0e-4}{m^3}$、$\rho=\qty{1.20}{kg/m^3}$、$W=+\qty{28}{J}$（縮むので正）、$Q=0$。
\phriyu
断熱なので $\Delta U=W$（⑤）。空気の内部エネルギーは温度だけで決まり（⑥）、温度を\qty{1}{K}上げるのに要る内部エネルギーの増えは、体積一定で温めるときの熱と同じです（⑦の $C_V$）。だから $\Delta U=mc_v\,\Delta T$ と書けます。
\phtateru
\[m=\rho V,\qquad mc_v\,\Delta T=W\]
\phtoku
\[m=1.20\times\num{2.0e-4}=\qty{2.4e-4}{kg}\ (\qty{0.24}{g})\]
\[\Delta T=\frac{28}{\num{2.4e-4}\times718}\approx\qty{1.6e2}{K}\]
\phtashikame
大きさ：$20+160=180\,℃$ になる計算で、熱-3 ⑥（例題6）で断熱変化の式から求める約182\,℃と合います。\\
仮定がくずれる：本物の空気入れでは、熱がすぐに筒へ逃げ、空気もタイヤへ流れ出すので、筒の先が「温かい」くらいです。\qty{0.24}{g}の空気の熱容量は約\qty{0.17}{J/K}しかなく、筒の金属（数十 g）にくらべてとても小さいからです。
\kotae 約\qty{1.6e2}{K}（約180\,℃になる計算）
\end{phreidai}
\begin{jibun}
空気\qty{0.30}{L}（密度\qty{1.20}{kg/m^3}）に、断熱で\qty{30}{J}の仕事をして押し縮めました。温度の上がりを求めましょう（$c_v=718$\,J/(kg·K)）。
\bkjans{$m=\qty{3.6e-4}{kg}$。$30\div(\num{3.6e-4}\times718)$ より、約\qty{1.2e2}{K}}
\end{jibun}

\label{ph:熱-2:ex10}
\begin{phreidai}[電気ストーブが部屋に入れるエネルギー]
\qty{1.0}{kW}の電気ストーブを1時間使いました。部屋に入るエネルギーは何 MJ ですか。もし、このエネルギーが全部、部屋の空気（\qty{50}{m^3}、\qty{60}{kg}、圧力一定の比熱 1.0\,kJ/(kg·K)）にたまったら、温度は何 K 上がりますか。
\phbamen
コンセントから入った電気のエネルギーが、ストーブを通って、部屋に入ります。
\phmushi
ストーブの熱容量を無視します（すぐに一定の温度になる）。部屋の空気の出入りを無視します。
\phzu
$P=\qty{1.0e3}{W}$、$t=\qty{3600}{s}$。空気 $m=\qty{60}{kg}$、$c=1000\,J/(kg·K)$。
\phriyu
ストーブを系にすると、電気の仕事 $W=Pt$ が入り、同じだけの熱 $Q$ が部屋へ出ていきます（$\Delta U\approx0$。④の「電気の仕事も、仕事」）。部屋の空気を系にすると、その熱を受けとって温まります（熱-1 ③）。
\phtateru
\[Q=Pt,\qquad mc\,\Delta T=Q\]
\phtoku
\[Q=\num{1.0e3}\times3600=\qty{3.6e6}{J}=\qty{3.6}{MJ},\qquad \Delta T=\frac{\num{3.6e6}}{60\times1000}=\qty{60}{K}\]
\phtashikame
仮定がくずれる：本物の部屋が\qty{60}{K}も温まらないのは、熱の大部分が、壁・床・家具に入ったり、外へ逃げたりするからです（熱-11）。\\
くらべる：電気の仕事は、全部熱になります（効率は1）。同じ電気で、外の空気から熱を運んでくるヒートポンプ（エアコン）なら、もっと多くの熱を部屋に入れられます（熱-5・熱-9）。
\kotae \qty{3.6}{MJ}。空気だけなら\qty{60}{K}上がる計算
\end{phreidai}
\begin{jibun}
\qty{600}{W}の電気ストーブを30分使うと、部屋に入るエネルギーは何 MJ ですか。
\bkjans{$600\times1800=\qty{1.08e6}{J}$ より、約\qty{1.1}{MJ}}
\end{jibun}

\label{ph:熱-2:ex11}
\begin{phreidai}[状態量の小さな変化はどちら？]
（1）$y\,\dd x+x\,\dd y$、（2）$-y\,\dd x+x\,\dd y$ のうち、ある関数の全微分（完全微分）になっているのはどちらですか。また、それぞれを、点 $(0,\ 0)$ から点 $(1,\ 1)$ まで、（A）$(1,\ 0)$ を通る道と、（B）$(0,\ 1)$ を通る道でたして、くらべましょう（力-4 例題13 と同じ形）。
\phbamen
完全微分なら、ある状態量の小さな変化なので、たした値は道によらないはずです。
\phmushi
数学の問題として扱います（$x$・$y$ は、たとえば $V$・$p$ のような状態を表す数と考える）。
\phzu
$A\,\dd x+B\,\dd y$ と書いて、（1）$A=y$、$B=x$、（2）$A=-y$、$B=x$。道Aは右へ進んでから上へ、道Bは上へ進んでから右へ。
\phriyu
⑦の見分け方 $\partial A/\partial y=\partial B/\partial x$ を使います。道ごとの和は、横の区間で $A\,\dd x$（$y$ は止める）、たての区間で $B\,\dd y$（$x$ は止める）です（8-7⑤）。
\phtateru
（1）$\partial A/\partial y=1$、$\partial B/\partial x=1$。（2）$\partial A/\partial y=-1$、$\partial B/\partial x=1$。
\phtoku
（1）は等しいので完全微分です。$f=xy$ とすると $\dd f=y\,\dd x+x\,\dd y$。道A：$\int_0^1 0\,\dd x+\int_0^1 1\,\dd y=1$、道B：$\int_0^1 0\,\dd y+\int_0^1 1\,\dd x=1$。\\
（2）は等しくないので、完全微分ではありません。道A：$\int_0^1 0\,\dd x+\int_0^1 1\,\dd y=1$、道B：$\int_0^1 0\,\dd y+\int_0^1(-1)\,\dd x=-1$。
\phtashikame
状態量で確かめる：（1）は $f(1,\ 1)-f(0,\ 0)=1-0=1$ で、道によりません。\\
熱力学では：（1）は $\dd U$ のなかま、（2）は $\delta Q$・$\delta W$ のなかまです。（2）を1周（Aで行ってBでもどる）たすと $1-(-1)=2$ で、0になりません（⑧のサイクルの正味の仕事と同じ考え）。
\kotae （1）が完全微分（どちらの道も1）。（2）は完全微分ではない（道Aで1、道Bで $-1$）
\end{phreidai}
\begin{jibun}
（1）$2xy\,\dd x+x^2\,\dd y$、（2）$y\,\dd x-x\,\dd y$ は、完全微分ですか。完全微分なら、もとの関数も答えましょう。
\bkjans{（1）$\partial(2xy)/\partial y=2x=\partial(x^2)/\partial x$ で完全微分。$f=x^2y$。（2）$\partial y/\partial y=1$、$\partial(-x)/\partial x=-1$ で、完全微分ではない}
\end{jibun}

\label{ph:熱-2:ex12}
\begin{phreidai}[長方形のサイクルの正味の仕事と熱]
⑧の図のサイクル A→B→C→D→A（圧力 200\,kPa と 100\,kPa、体積 1\,L と 3\,L の長方形を時計まわり）で、1周の間に気体が外にする正味の仕事と、受けとる正味の熱を求めましょう。
\phbamen
気体が、ふくらんで外へ仕事をし、縮んで仕事をされて、元にもどります。
\phmushi
どの変化も準静的とします。
\phzu
A→B：200\,kPa で 1→3\,L。B→C：3\,L のまま。C→D：100\,kPa で 3→1\,L。D→A：1\,L のまま。kPa × L ＝ J。
\phriyu
区間ごとに $W'=p\,\Delta V$（③）をたします。体積が変わらない区間は0です。1周で $\oint\dd U=0$ なので、正味の熱は正味の仕事に等しくなります（⑧）。
\phtateru
\[W_{\text{出}}-W_{\text{入}}=W'_{\mathrm{AB}}+W'_{\mathrm{BC}}+W'_{\mathrm{CD}}+W'_{\mathrm{DA}},\qquad Q_{\text{入}}-Q_{\text{出}}=W_{\text{出}}-W_{\text{入}}\]
\phtoku
\[W_{\text{出}}-W_{\text{入}}=200\times(3-1)+0+100\times(1-3)+0=400-200=\qty{200}{J}\]
\[Q_{\text{入}}-Q_{\text{出}}=\qty{200}{J}\]
\phtashikame
面積で：長方形の面積 $(200-100)\times(3-1)=\qty{200}{J}$ と同じです。\\
向き：反時計まわりに回すと、$W_{\text{出}}-W_{\text{入}}=-\qty{200}{J}$ になり、外から仕事をされて、同じだけの熱を出します（冷凍機の向き。熱-5）。
\kotae 正味の仕事 \qty{200}{J}、正味の熱 \qty{200}{J}
\end{phreidai}
\begin{jibun}
圧力 400\,kPa と 100\,kPa、体積 1\,L と 2\,L の長方形を時計まわりに1周するサイクルで、気体が外にする正味の仕事を求めましょう。
\bkjans{$(400-100)\times(2-1)=\qty{300}{J}$}
\end{jibun}

%% ======================== まちがえやすいところ ========================
\Needspace{11\baselineskip}\bksection{tsumazuki}{まちがえやすいところをチェック}
\begin{tsumazuki}{$W$ は、するほう？　されるほう？}
\bkng{気体がふくらんで外へ\qty{100}{J}の仕事をしたので、$\Delta U=Q+100$。}
\bkok{この本の $W$ は「される」仕事。ふくらむときは $W=-\qty{100}{J}$ で、$\Delta U=Q-100$。}
本や問題によって、$W$ の向きの決め方が違います。迷ったら、矢印の図をかいて、入る量をたし、出る量をひきます（④）。
\end{tsumazuki}

\begin{tsumazuki}{熱を加えたら、必ず温度が上がる？}
\bkng{気体に熱を加えたので、内部エネルギーが増え、温度が上がる。}
\bkok{熱を加えても、それ以上の仕事を外へすれば、内部エネルギーは減る。}
$\Delta U=Q-W'$ なので、$W'>Q$ なら $\Delta U<0$ です。熱と温度は別のものです（熱-1 ①）。
\end{tsumazuki}

\begin{tsumazuki}{仕事は、はじめとおわりの状態で決まる？}
\bkng{はじめとおわりの圧力と体積がわかれば、気体がした仕事は決まる。}
\bkok{仕事と熱は、道すじで変わる。決まるのは状態量の変化 $\Delta U$ だけ。}
8-7⑩の3つの道で、仕事は 400\,J・200\,J・300\,J と違いました（③⑦）。
\end{tsumazuki}

\begin{tsumazuki}{断熱なら、温度は変わらない？}
\bkng{熱が出入りしない（断熱）ので、温度は変わらない。}
\bkok{断熱でも、仕事をされれば温度は上がる。$\Delta U=W$。}
断熱は「$Q=0$」という意味で、「温度一定」ではありません。温度一定のまま押し縮めるには、熱を外へ逃がす必要があります（⑤）。
\end{tsumazuki}

%% ======================== 工学ではこう使う ========================
\Needspace{18\baselineskip}\bksection{phgear}{工学ではこう使う}
\begin{phkougaku}[ブレーキの熱を、どこへ逃がすか]
車のブレーキは、運動エネルギーを内部エネルギー（熱）に変えて止める装置です（例題6）。熱がたまって円板やパッドが高温になりすぎると、摩擦係数が下がってブレーキが効きにくくなります（フェード現象）。そこで、円板の中に風の通る羽根を作った「ベンチレーテッドディスク」で、空気へ熱を逃がしやすくしています（対流で熱を運ぶしくみは熱-12）。長い下り坂では、エンジンブレーキ（シフトダウン）を使って、ブレーキに入る熱を減らします。電車や電気自動車は、運動エネルギーの一部を電気にもどしてためます（回生ブレーキ、力-4）。熱になったエネルギーを、すべて仕事にもどすことはできないからです（熱-5）。

\begin{phketa}[熱と仕事の大きさ]
\phketaline{\qty{1}{kg}の物を\qty{1}{m}持ち上げる仕事}{\qty{9.8}{J}}
\phketaline{水\qty{1}{kg}を\qty{1}{K}上げる熱}{\qty{4.2}{kJ}（約430倍）}
\phketaline{車\qty{1200}{kg}・\qty{20}{m/s}の運動エネルギー}{\qty{0.24}{MJ}}
\phketaline{電気ストーブ\qty{1}{kW}を1時間}{\qty{3.6}{MJ}（\qty{1}{kWh}）}
\phboxhead{答えが桁で合っているか}
力学の仕事（数 J〜数百 kJ）は、水を温める熱にすると小さい量です。「人が数十 m の階段を上ったら、体温が何十 K も上がる」のような答えが出たら、J と kJ の取り違いを疑います（\qty{60}{kg}の人が\qty{30}{m}上る仕事は約\qty{18}{kJ}で、体温なら0.1\,K 分にもなりません）。
\end{phketa}
\end{phkougaku}
\phstation{熱力学,内燃機関,ブレーキと車両,化学工学}{内部エネルギー、第一法則、閉じた系、準静的過程、絶対仕事、断熱変化、自由膨張、状態量、サイクル}

\begin{column}{ジュールと「エネルギー保存」}
ジェームズ・プレスコット・ジュール（1818〜1889年）は、イギリスのビール工場の家に生まれ、自分の家で実験を続けた人でした。電流の熱（ジュール熱）の研究から出発して、「熱は仕事や電気から作ることができ、換算の比はいつも同じ」ことを、羽根車などの実験で確かめました。はじめのころの発表は、あまり信じてもらえませんでしたが、1847年の学会で若いウィリアム・トムソン（のちのケルビン）が興味をもち、やがて広く認められました。同じころ、ドイツのマイヤー（1842年）やヘルムホルツ（1847年）も、別の道すじから「エネルギーは形を変えても保存する」と考えました。エネルギーの単位 J（ジュール）は、ジュールの名前からとられています。
\end{column}

\begin{column}{燃料なしで動き続ける機械は、なぜつくれない？}
外から何も入れないのに、仕事を出し続ける機械を\textbf{永久機関}といいます。昔から、重りや水車や磁石を組み合わせた、たくさんの「永久機関」が考えられてきました。でも、どれも、しばらくすると止まってしまいました。1775年には、フランスの科学アカデミーが、永久機関の提案はもう審査しない、と決めています。

第一法則で考えると、理由がはっきりします。機械が1周して元の状態にもどるなら、$\oint\dd U=0$ なので、外にした正味の仕事は、受けとった正味の熱に等しくなります（⑧）。熱も仕事も外から入れないなら、出てくる正味の仕事は0です。エネルギーを入れずに仕事だけを取り出す機械（第1種永久機関）は、第一法則にそむくので、つくれないのです。では、「まわりの空気や海の熱を受けとって、それを全部仕事に変え続ける機械」ならどうでしょう。これは第一法則にはそむきません。それでもつくれないわけを、熱-5 で学びます。
\end{column}

%% ======================== 章末問題 ========================
\bkendproblems
\bkproblemgroup{基本}{2}
\bkproblem 質量\qty{5.0}{kg}の重りが\qty{1.0}{m}落ちるのを50回繰り返して、断熱した入れ物の中の水\qty{0.50}{kg}をかき回しました。水の温度の上がりを求めましょう。
\bkpans{$W=50\times5.0\times9.8\times1.0=\qty{2450}{J}$。$2450\div(0.50\times4200)$ より、約\qty{1.2}{K}}
\bkproblem 気体の内部エネルギーの変化 $\Delta U$ を求めましょう。（1）\qty{300}{J}の熱を受けとり、外へ\qty{120}{J}の仕事をした。（2）断熱で、外から\qty{50}{J}の仕事をされた。（3）体積を変えずに、\qty{80}{J}の熱を外へ出した。
\bkpans{（1）\nobreak $300+(-120)=\qty{180}{J}$。（2）\nobreak $+\qty{50}{J}$。（3）\nobreak $-\qty{80}{J}$}
\bkproblem 圧力\qty{150}{kPa}のまま、気体の体積が\qty{2.0}{L}から\qty{5.0}{L}にふくらみました。（1）気体がした仕事 $W'$ と、された仕事 $W$ を求めましょう。（2）この間に $\Delta U=\qty{1125}{J}$ だったとき、気体が受けとった熱を求めましょう。
\bkpans{（1）\nobreak $W'=\qty{450}{J}$、$W=-\qty{450}{J}$。（2）\nobreak $Q=\Delta U-W=1125+450$ より、約\qty{1.6e3}{J}（1575）}
\bkproblem 人と自転車で\qty{80}{kg}が、\qty{8.0}{m/s}からブレーキで止まりました。ブレーキで熱になったエネルギーを求めましょう（道は水平、空気の抵抗は無視）。
\bkpans{$\tfrac{1}{2}\times80\times8.0^2=\qty{2560}{J}$ より、約\qty{2.6}{kJ}}
\bkproblem 電気ポットの水\qty{1.0}{kg}を\qty{1.0}{kW}のヒーターで\qty{180}{s}温める間に、壁から\qty{18}{kJ}の熱が逃げました。水とヒーターを系とし、（1）系がされた電気の仕事、（2）$\Delta U$、（3）水の温度の上がりを求めましょう（ヒーターの熱容量は無視）。
\bkpans{（1）\nobreak $W=\qty{180}{kJ}$。（2）\nobreak $\Delta U=Q+W=-18+180=\qty{162}{kJ}$。（3）\nobreak $162\div4.2$ より、約\qty{39}{K}}
\bkproblem 気体の状態を、$p$-$V$ 図の上で、（1\,L, 100\,kPa）→（3\,L, 100\,kPa）→（3\,L, 300\,kPa）→（1\,L, 100\,kPa）と、3本の線分にそって1周させました。1周で気体が外にする正味の仕事 $W_{\text{出}}-W_{\text{入}}$ と、受けとる正味の熱を、回る向きに注意して求めましょう。
\bkpans{1つめ $+\qty{200}{J}$、2つめ 0、3つめ（平均の圧力 200\,kPa で 2\,L 縮む）$-\qty{400}{J}$。$W_{\text{出}}-W_{\text{入}}=-\qty{200}{J}$（反時計まわり）。正味の熱も $-\qty{200}{J}$（\qty{200}{J}出す）}
\bkproblemgroup{チャレンジ}{4}
\bkproblem 断熱した壁の容器が、しきりで同じ体積\qty{10}{L}の2つの部屋に分かれています。左には\qty{300}{K}の理想気体、右は真空です。しきりを外すと、気体は容器全体に広がりました。（1）この変化の $Q$・$W$・$\Delta U$ を答えましょう。（2）広がったあとの温度を求めましょう。（3）圧力は、はじめの何倍になりますか。
\bkpans{（1）\nobreak $Q=0$（断熱）、$W=0$（真空へ広がる）、$\Delta U=0$。（2）\nobreak 理想気体の $U$ は温度だけで決まるので、\qty{300}{K}のまま。（3）\nobreak 温度が同じで体積が2倍なので、$pV$ 一定（3-7⑥）から $\tfrac{1}{2}$ 倍}
\bkproblem $(3x^2y+2)\,\dd x+(x^3+2y)\,\dd y$ が完全微分であることを確かめ、もとの関数 $f(x,\ y)$ を求めましょう。また、点 $(0,\ 0)$ から点 $(1,\ 2)$ まで、この量をたした値を求めましょう。
\bkpans{$\partial(3x^2y+2)/\partial y=3x^2=\partial(x^3+2y)/\partial x$ で完全微分。$f=x^3y+2x+y^2$（＋定数）。たした値は $f(1,\ 2)-f(0,\ 0)=2+2+4=8$（道によらない）}
\bkproblem 例題9の空気入れで、仕事は\qty{28}{J}でしたが、空気は\qty{60}{K}しか温まりませんでした（\qty{0.24}{g}、$c_v=718$\,J/(kg·K)）。筒へ逃げた熱を求めましょう。
\bkpans{$\Delta U=\num{2.4e-4}\times718\times60\approx\qty{10}{J}$。$Q=\Delta U-W=10.3-28$ より、約\qty{18}{J}が逃げた（$Q\approx-\qty{18}{J}$）}

\begin{furikaeri}
\begin{checklist}
\item 熱と仕事が同じエネルギーであることを、実験から説明できる。\bkfurref{ph:熱-2:sec1}
\item 第一法則 $\Delta U=Q+W$ で、エネルギーの出入りを数えられる。\bkfurref{ph:熱-2:sec4}
\item $p$-$V$ 図の面積から、気体がする仕事を求められる。\bkfurref{ph:熱-2:sec3}
\item 状態量と、道すじで変わる量を、全微分で見分けられる。\bkfurref{ph:熱-2:sec7}
\end{checklist}
\end{furikaeri}
\bkfukidashi{mina}{熱と仕事は、エネルギーの2つの入り口なんだね。空気入れの先が熱いのは、押した仕事が、空気の内部エネルギーになったからだ。}
